\documentclass[a4paper,12pt,twoside]{memoir}

% Castellano
\usepackage[spanish,es-tabla]{babel}
\selectlanguage{spanish}
\usepackage[utf8]{inputenc}
\usepackage[T1]{fontenc}
\usepackage{lmodern} % Scalable font
\usepackage{microtype}
\usepackage{placeins}

\RequirePackage{booktabs}
\RequirePackage[table]{xcolor}
\RequirePackage{xtab}
\RequirePackage{multirow}

% Links
\usepackage[colorlinks]{hyperref}
\hypersetup{
	allcolors = {red}
}

% Ecuaciones
\usepackage{amsmath}

% Rutas de fichero / paquete
\newcommand{\ruta}[1]{{\sffamily #1}}

% Párrafos
\nonzeroparskip

% Huérfanas y viudas
\widowpenalty100000
\clubpenalty100000

% Evitar solapes en el header
\nouppercaseheads

% Imagenes
\usepackage{graphicx}
\newcommand{\imagen}[2]{
	\begin{figure}[!h]
		\centering
		\includegraphics[width=0.95\textwidth]{#1}
		\caption{#2}\label{fig:#1}
	\end{figure}
	\FloatBarrier
}

\newcommand{\imagenSinEspacio}[2]{
	\begin{figure}
		\centering
		\includegraphics{#1}
		\caption{#2}\label{fig:#1}
	\end{figure}
	\FloatBarrier
}

\newcommand{\miniimagen}[2]{
	\begin{figure}[!h]
		\centering
		\includegraphics[width=0.5\textwidth]{#1}
		\caption{#2}\label{fig:#1}
	\end{figure}
	\FloatBarrier
}

\newcommand{\imagenflotante}[2]{
	\begin{figure}%[!h]
		\centering
		\includegraphics[width=0.9\textwidth]{#1}
		\caption{#2}\label{fig:#1}
	\end{figure}
}


% El comando \figura nos permite insertar figuras comodamente, y utilizando
% siempre el mismo formato. Los parametros son:
% 1 -> Porcentaje del ancho de página que ocupará la figura (de 0 a 1)
% 2 --> Fichero de la imagen
% 3 --> Texto a pie de imagen
% 4 --> Etiqueta (label) para referencias
% 5 --> Opciones que queramos pasarle al \includegraphics
% 6 --> Opciones de posicionamiento a pasarle a \begin{figure}
\newcommand{\figuraConPosicion}[6]{%
  \setlength{\anchoFloat}{#1\textwidth}%
  \addtolength{\anchoFloat}{-4\fboxsep}%
  \setlength{\anchoFigura}{\anchoFloat}%
  \begin{figure}[#6]
    \begin{center}%
      \Ovalbox{%
        \begin{minipage}{\anchoFloat}%
          \begin{center}%
            \includegraphics[width=\anchoFigura,#5]{#2}%
            \caption{#3}%
            \label{#4}%
          \end{center}%
        \end{minipage}
      }%
    \end{center}%
  \end{figure}%
}

%
% Comando para incluir imágenes en formato apaisado (sin marco).
\newcommand{\figuraApaisadaSinMarco}[5]{%
  \begin{figure}%
    \begin{center}%
    \includegraphics[angle=90,height=#1\textheight,#5]{#2}%
    \caption{#3}%
    \label{#4}%
    \end{center}%
  \end{figure}%
}
% Para las tablas
\newcommand{\otoprule}{\midrule [\heavyrulewidth]}
%
% Nuevo comando para tablas pequeñas (menos de una página).
\newcommand{\tablaSmall}[5]{%
 \begin{table}
  \begin{center}
   \rowcolors {2}{gray!35}{}
   \begin{tabular}{#2}
    \toprule
    #4
    \otoprule
    #5
    \bottomrule
   \end{tabular}
   \caption{#1}
   \label{tabla:#3}
  \end{center}
 \end{table}
}

%
% Nuevo comando para tablas pequeñas (menos de una página).
\newcommand{\tablaSmallSinColores}[5]{%
 \begin{table}[H]
  \begin{center}
   \begin{tabular}{#2}
    \toprule
    #4
    \otoprule
    #5
    \bottomrule
   \end{tabular}
   \caption{#1}
   \label{tabla:#3}
  \end{center}
 \end{table}
}

\newcommand{\tablaApaisadaSmall}[5]{%
\begin{landscape}
  \begin{table}
   \begin{center}
    \rowcolors {2}{gray!35}{}
    \begin{tabular}{#2}
     \toprule
     #4
     \otoprule
     #5
     \bottomrule
    \end{tabular}
    \caption{#1}
    \label{tabla:#3}
   \end{center}
  \end{table}
\end{landscape}
}

%
% Nuevo comando para tablas grandes con cabecera y filas alternas coloreadas en gris.
\newcommand{\tabla}[6]{%
  \begin{center}
    \tablefirsthead{
      \toprule
      #5
      \otoprule
    }
    \tablehead{
      \multicolumn{#3}{l}{\small\sl continúa desde la página anterior}\\
      \toprule
      #5
      \otoprule
    }
    \tabletail{
      \hline
      \multicolumn{#3}{r}{\small\sl continúa en la página siguiente}\\
    }
    \tablelasttail{
      \hline
    }
    \bottomcaption{#1}
    \rowcolors {2}{gray!35}{}
    \begin{xtabular}{#2}
      #6
      \bottomrule
    \end{xtabular}
    \label{tabla:#4}
  \end{center}
}

%
% Nuevo comando para tablas grandes con cabecera.
\newcommand{\tablaSinColores}[6]{%
  \begin{center}
    \tablefirsthead{
      \toprule
      #5
      \otoprule
    }
    \tablehead{
      \multicolumn{#3}{l}{\small\sl continúa desde la página anterior}\\
      \toprule
      #5
      \otoprule
    }
    \tabletail{
      \hline
      \multicolumn{#3}{r}{\small\sl continúa en la página siguiente}\\
    }
    \tablelasttail{
      \hline
    }
    \bottomcaption{#1}
    \begin{xtabular}{#2}
      #6
      \bottomrule
    \end{xtabular}
    \label{tabla:#4}
  \end{center}
}

%
% Nuevo comando para tablas grandes sin cabecera.
\newcommand{\tablaSinCabecera}[5]{%
  \begin{center}
    \tablefirsthead{
      \toprule
    }
    \tablehead{
      \multicolumn{#3}{l}{\small\sl continúa desde la página anterior}\\
      \hline
    }
    \tabletail{
      \hline
      \multicolumn{#3}{r}{\small\sl continúa en la página siguiente}\\
    }
    \tablelasttail{
      \hline
    }
    \bottomcaption{#1}
  \begin{xtabular}{#2}
    #5
   \bottomrule
  \end{xtabular}
  \label{tabla:#4}
  \end{center}
}

\definecolor{cgoLight}{HTML}{EEEEEE}
\definecolor{cgoExtralight}{HTML}{FFFFFF}

%
% Nuevo comando para tablas grandes sin cabecera.
\newcommand{\tablaSinCabeceraConBandas}[5]{%
  \begin{center}
    \tablefirsthead{
      \toprule
    }
    \tablehead{
      \multicolumn{#3}{l}{\small\sl continúa desde la página anterior}\\
      \hline
    }
    \tabletail{
      \hline
      \multicolumn{#3}{r}{\small\sl continúa en la página siguiente}\\
    }
    \tablelasttail{
      \hline
    }
    \bottomcaption{#1}
    \rowcolors[]{1}{cgoExtralight}{cgoLight}

  \begin{xtabular}{#2}
    #5
   \bottomrule
  \end{xtabular}
  \label{tabla:#4}
  \end{center}
}


\graphicspath{ {./img/} }

% Capítulos
\chapterstyle{bianchi}
\newcommand{\capitulo}[2]{
	\setcounter{chapter}{#1}
	\setcounter{section}{0}
	\chapter*{#2}
	\addcontentsline{toc}{chapter}{#1. #2}
	\markboth{#2}{#2}
}

% Apéndices
\renewcommand{\appendixname}{Apéndice}
\renewcommand*\cftappendixname{\appendixname}

\newcommand{\apendice}[1]{
	%\renewcommand{\thechapter}{A}
	\chapter{#1}
}

\renewcommand*\cftappendixname{\appendixname\ }

% Formato de portada
\makeatletter
\usepackage{xcolor}
\newcommand{\tutor}[1]{\def\@tutor{#1}}
\newcommand{\course}[1]{\def\@course{#1}}
\definecolor{cpardoBox}{HTML}{E6E6FF}
\def\maketitle{
  \null
  \thispagestyle{empty}
  % Cabecera ----------------
\begin{center}%
	{\noindent\Huge Universidades de Burgos, León y Valladolid}\vspace{.5cm}%
	
	{\noindent\Large Máster universitario}\vspace{.5cm}%
	
	{\noindent\Huge \textbf{Inteligencia de Negocio y Big~Data en Entornos Seguros}}\vspace{.5cm}%
\end{center}%

\begin{center}%
	\includegraphics[height=3cm]{img/escudoUBU} \hspace{1cm}
	\includegraphics[height=3cm]{img/escudoUVA} \hspace{1cm}
	\includegraphics[height=3cm]{img/escudoULE} \vspace{1cm}%
\end{center}%

  \vfill
  % Título proyecto y escudo informática ----------------
  \colorbox{cpardoBox}{%
    \begin{minipage}{.9\textwidth}
      \vspace{.1cm}\Large
      \begin{center}
      \textbf{TFM del Máster Inteligencia de Negocio y Big Data en Entornos Seguros}\vspace{.6cm}\\
      \textbf{\large\@title{}}
      \end{center}
      \vspace{.2cm}
    \end{minipage}

  }%
  \hfill
  \vfill
  % Datos de alumno, curso y tutores ------------------
  \begin{center}%
  {%
    \noindent\large
    Presentado por \@author{}\\ 
    en Universidad de Burgos --- \@date{}\\
    Tutor: \@tutor{}\\
  }%
  \end{center}%
  \null
  \cleardoublepage
  }
\makeatother

\newcommand{\nombre}{Mario Varona Bueno} %%% cambio de comando

% Datos de portada
\title{Sistema de recomendación de municipios españoles en línea basado en las preferencias del usuario y en la similitud entre municipios.}
\author{\nombre}
\tutor{Álvar Arnaiz González}
\date{15 de enero de 2023}

\begin{document}

\maketitle

\newpage\null\thispagestyle{empty}\newpage

%%%%%%%%%%%%%%%%%%%%%%%%%%%%%%%%%%%%%%%%%%%%%%%%%%%%%%%%%%%%%%%%%%%%%%%%%%%%%%%%%%%%%%%%
\thispagestyle{empty}


\noindent
\begin{center}%
	{\noindent\Huge Universidades de Burgos, León y Valladolid}\vspace{.5cm}%
	
\begin{center}%
	\includegraphics[height=3cm]{img/escudoUBU} \hspace{1cm}
	\includegraphics[height=3cm]{img/escudoUVA} \hspace{1cm}
	\includegraphics[height=3cm]{img/escudoULE} \vspace{1cm}%
\end{center}%

	{\noindent\Large \textbf{Máster universitario en Inteligencia de Negocio y Big~Data en Entornos Seguros}}\vspace{.5cm}%
\end{center}%

\noindent D. Álvar Arnaiz González, profesor del departamento de Ingeniería Informática, área de Lenguajes y Sistemas Informáticos.

\noindent Expone:

\noindent Que el alumno D. \nombre, con DNI 70918530P, ha realizado el Trabajo final de Máster en Inteligencia de Negocio y Big Data en Entornos Seguros 
          titulado ``Sistema de recomendación de municipios españoles en línea basado en las preferencias del usuario y en la similitud entre municipios''. 

\noindent Y que dicho trabajo ha sido realizado por el alumno bajo la dirección del que suscribe, en virtud de lo cual se autoriza su presentación y defensa.

\begin{center} %\large
En Burgos, {\large 15 de enero de 2023}
\end{center}

\vfill\vfill\vfill

% Author and supervisor
% \begin{minipage}{0.45\textwidth}
% \begin{flushleft} %\large
% Vº. Bº. del Tutor:\\[2cm]
% 
% \end{flushleft}
% \end{minipage}
% \hfill
% \begin{minipage}{0.45\textwidth}
% \begin{flushleft} %\large
% Vº. Bº. del co-tutor:\\[2cm]
% 
% \end{flushleft}
% \end{minipage}
% \hfill

\vfill

% para casos con solo un tutor comentar lo anterior
% y descomentar lo siguiente
Vº. Bº. del Tutor:\\[2cm]
Álvar Arnaiz González

\newpage\null\thispagestyle{empty}\newpage

\frontmatter

% Abstract en castellano
\renewcommand*\abstractname{Resumen}
\begin{abstract}
En España existen 8.131 municipios, compuestos por más de 60.000 entidades singulares de población. Sin embargo, la gran mayoría de ellos son totalmente desconocidos incluso para personas de la provincia.

La mejora de las telecomunicaciones, infraestructuras, desarrollos turísticos e incremento del trabajo remoto, entre otros muchos factores, permiten considerar nuevos destinos como lugares de ocio o asentamiento; y la minería de datos y las técnicas de procesamiento de grandes cantidades de información permiten ofrecer recomendaciones personalizadas para cada usuario en función de sus gustos.

En el presente trabajo de fin de máster se abordará el diseño y creación de un sistema de recomendación en línea para todos los municipios españoles, que ofrecerá al usuario aquellos que mejor se adapten a sus gustos y más se parezcan a otros cuya preferencia ha manifestado anteriormente, entre otras características.

Primeramente, se expondrá el problema a resolver y los diferentes enfoques existentes. Posteriormente, se presentará la solución elegida junto con los retos encontrados y superados. Finalmente, se realizará una evaluación del funcionamiento del sistema, con la intención de comprobar su utilidad real y de medir el grado de consecución de los objetivos propuestos.
\end{abstract}

\renewcommand*\abstractname{Descriptores}
\begin{abstract}
Sistema de recomendación en línea, municipios, España, ciudades, pueblos, sistema de recomendación basado en contenido, sistema de recomendación con filtro colaborativo. 
\end{abstract}

\clearpage

% Abstract en inglés
\renewcommand*\abstractname{Abstract}
\begin{abstract}
There are 8,131 municipalities in Spain, which are compound by more than 64,000 singular entities of population. However, the vast majority of them are totally unknown even by local people.

The improvement of telecommunications, infrastructures, tourist developments and rise of remote working, among many other factors, allow considering new destinations as places of leisure or settlement; and data mining and processing techniques of large amounts of information allow offering personalized recommendations for each user according to their tastes.

This master's thesis will address the design and creation of an online recommendation system for all Spanish municipalities, which will offer to users those places that suit best their tastes and are most similar to others whose preference has been previously expressed, among other characteristics.

Firstly, the problem to be solved and the different existing approaches will be presented. Next, the chosen solution will be introduced together with the encountered and overcome challenges. Finally, an evaluation of the system will be carried out, in order to verify its actual utility and to measure the achievement of the proposed goals.
\end{abstract}

\renewcommand*\abstractname{Keywords}
\begin{abstract}
Online recommendation system, municipalities, Spain, cities, towns, villages, content-based recommendation system, collaborative filter recommendation system.
\end{abstract}

\clearpage

% Indices
\tableofcontents

\clearpage

\listoffigures

\clearpage

\mainmatter

\addcontentsline{toc}{part}{Memoria}
\part*{Memoria}

\capitulo{1}{Introducción}

Los sistemas de recomendación se han convertido en una de las aplicaciones más usadas dentro del abanico de posibilidades que ofrecen las tecnologías que explotan \textit{Big Data}; es decir, que se aprovechan de una gran cantidad de datos para funcionar. Encontramos estos sistemas de recomendación en redes sociales, portales audiovisuales o cualquier otra plataforma de consumo de contenido –como libros o noticias–, donde nos pueda interesar conocer más sobre alguna temática por la que hemos manifestado un interés previamente. Estos sistemas están, por tanto, especialmente indicados para casos de uso en los que la oferta posible supera con creces la facilidad natural con la que el contenido sería descubierto por los interesados en ausencia del mecanismo de sugerencias automáticas. Si se implementan correctamente, estos mecanismos se convierten en componentes del sistema en los que el usuario confía a la hora de descubrir nuevos grupos musicales, información relacionada o, por qué no, destinos vacacionales o permanentes, atribuyéndoles un argumento de autoridad que en el mundo analógico es habitualmente otorgado a personas de confianza para el usuario, o que considera expertos en la materia.

Partiendo de esa base, y aplicando todos los contenidos vistos en el Máster para el que esta memoria constituye el entregable principal del trabajo final, se ha deseado construir un sistema de recomendación con diversas características que se irán presentando a lo largo de la misma. Pese a que se hace un uso extenso de los conocimientos adquiridos en las asignaturas ``Aprendizaje No Supervisado'' y ``Técnicas de Aprendizaje Automático Escalables'', se ha querido aprovechar el mayor número posible de competencias adquiridas en el resto de materias, con especial mención a las procedentes de ``Modelos de Programación para el Big Data'', ``Infraestructura para el Big Data'', ``Visualización de Datos'', ``Derecho en Seguridad de Datos'' y ``Fundamentos de Ciberseguridad''. En cualquier caso, las aportaciones de cada asignatura se señalarán convenientemente a lo largo de esta memoria.

Con la intención de aunar estos conocimientos en un trabajo teórico-práctico cuyo resultado pueda ser de utilidad, y dentro de la función de transferencia y aplicación de conocimiento que el alumno considera que debe exigirse al término de unos estudios universitarios, se propone como idea el diseño y creación de un sistema de recomendación de municipios de toda España, elegida por los motivos que a continuación se exponen:

\begin{itemize}
    \item \textbf{Necesidad de tratar con una gran cantidad de datos.} Pese a que en un primer momento se barajó la idea de acotar el ámbito a municipios de Castilla y León (2.248 frente a los 8.131 nacionales), y aprovechando la gran cantidad de datos abiertos en formato reusable que esta comunidad ofrece y que facilitaría en gran medida las labores de extracción, transformación, limpieza e ingestión de datos en el sistema, se descartó por diversos motivos: En primer lugar, las labores de recogida y limpieza de datos son tan importantes en un proyecto de Big Data como cualquier otra fase, consumiendo habitualmente entre un 50 \% y un 80 \% de los recursos del proyecto \cite{data_prep}. Esta regla asomaba intuitivamente en asignaturas que requerían algún tipo de ingestión de datos, como ``Arquitecturas Big Data'' y ``Almacenamiento Escalable'', pero no ha sido hasta la elaboración de este trabajo –en el que los datos no se suministraban como parte de un enunciado– cuando el alumno la ha podido comprobar empíricamente. En su caso, y dado que ha procurado usar fuentes de datos abiertas en todo lo posible (de ámbito nacional siempre que han estado disponibles, con la intención de evitar manejar fuentes de 19 comunidades y ciudades autónomas distintas), alrededor de un tercio del tiempo de desarrollo total del trabajo han sido invertidos en esta fase, lo que está estrechamente relacionado con el siguiente punto.
    \item \textbf{Deseo de aplicar técnicas de minería y extracción de datos.} Si bien es cierto que en el máster cursado se ha profundizado en diferentes técnicas para mantener una gran cantidad de datos en infraestructuras adecuadas (que escalen horizontalmente y permitan lecturas y escrituras rápidas), en maneras de obtener conclusiones sobre ellos gracias a técnicas estadísticas y de inteligencia artificial, en formas de mantenerlos seguros a nivel técnico y jurídico, y de visualizarlos cómodamente; las limitaciones lógicas de temario y tiempo han impedido incidir en técnicas de extracción automática de datos. Ha sido gracias a otros recursos aprendidos de forma autodidacta, como el alumno ha conocido técnicas de extracción automática de datos en Internet, cuya puesta en práctica se perfilaba casi obligada en el trabajo final de estos estudios.
    \item \textbf{Intención de crear un sistema de recomendación basado en contenido y en la retroalimentación de los usuarios.} El alumno deseaba poner en práctica sus conocimientos de los dos principales tipos de recomendadores: los basados en contenido y los que utilizan un filtro colaborativo empleando las preferencias de otros usuarios. A lo largo de esta memoria se abordarán ambos y su resultado final.
    \item \textbf{Voluntad de aportar una solución para descubrir ciudades y pueblos españoles.} En numerosas ocasiones se presenta el dilema de qué sitio elegir como destino vacacional o residencial. Hasta ahora, el lugar donde se encontrara trabajo condicionaba enormemente dónde se podría establecer una persona, y los destinos turísticos clásicos ocupaban la mayor parte de la oferta. En la actualidad, la masificación de los mismos ha traído consigo un deseo de explorar nuevas zonas, la mejora de las comunicaciones e infraestructuras ha provocado un resurgimiento en el interés por las zonas rurales, y el coste de vida asociado a las grandes urbes –destino de trabajo obligado para mucha gente– ha provocado la añoranza de una vida más cómoda, fácil de conseguir en una ciudad mediana o pequeña \cite{elpais}, \cite{terrenos}. Por su parte, el incremento del trabajo remoto desde la situación socio-sanitaria derivada de la crisis de la CoViD-19, ha materializado estos factores en un éxodo urbano experimentado en un porcentaje de población lo suficientemente significativo como para ser tenido en cuenta \cite{innovando}, \cite{20mins}. En nuestro país, el caso paradigmático es el de la ciudad de Madrid, que hasta ahora ha aglutinado a personas emigrantes del interior peninsular en busca de oportunidades laborales, pero no es el único. En el sur y archipiélagos peninsulares se ha observado el caso de trabajadores remotos de todo el mundo que han acudido a ellos para realizar sus funciones desde allí, y la tendencia es claramente alcista \cite{expansion}. Desde el ámbito mediático y político se puede apreciar como frecuentemente se iguala el concepto de ``España vaciada'' con territorios muy poco habitados o despoblados \cite{muycomputerpro}, olvidando que, para muchas personas, las limitaciones en municipios así son muy grandes como para plantearse el cambio desde una gran ciudad. Sin embargo, para el autor, esta España despoblada incluye inexorablemente a ciudades, capitales de provincia en muchos casos, que en las últimas décadas han visto disminuir su población drásticamente hacia los sumideros de las grandes áreas urbanas nacionales. Es el caso de la comunidad de las tres universidades de este máster, y es un motivo más para que el alumno no quisiera excluir a ciudades tan válidas para visitar o vivir como Burgos, León, Valladolid o Zamora, por poner algunos ejemplos. Además, y como se detallará en secciones posteriores, el trabajo con las capitales de provincia ha permitido aportar datos muy relevantes, que complementan la información de los municipios cuando, por su tamaño, no es posible disponer de esos datos.
    \item \textbf{Aprovechamiento de experiencia previa con datos abiertos.} El alumno ha querido aprovechar su conocimiento y trabajo previo con diversos catálogos de datos abiertos, como los de la Junta de Castilla y León, los del Gobierno de España o los de la Unión Europea. A mayores, ha querido experimentar con catálogos de datos procedentes de iniciativas privadas o extracción de datos existentes en Internet.
    \item \textbf{Materialización de una idea premiada previamente.} La idea de este trabajo fue una de las seis premiadas el pasado curso en el marco del Programa de Prototipos TCUE (Transferencia de Conocimiento Universidad-Empresa) de la Fundación Universidad de Burgos, dirigido a potenciar la creación de prototipos a partir de trabajos de fin de grado, máster o tesis doctorales que puedan aportar algo novedoso al mercado \cite{ubu}. Por otro lado, también fue premiada con el segundo premio en la categoría ``Ideas'' del V Concurso de Datos Abiertos de la Junta de Castilla y León \cite{cyl}, por lo que al recibir esta validación externa, se consideró que se debía intentar llevar a cabo. 
    \item \textbf{Deseo de aportar una solución novedosa.} El alumno considera que, además de demostrar y aplicar los conocimientos adquiridos a lo largo de los estudios universitarios, los trabajos finales de grado y máster también deben poder ser publicables; en el sentido de poder devolver a la sociedad parte de lo adquirido gracias a una formación pública de calidad que ella misma ha sufragado. Por ello, se ha diseñado también una forma de distribución de los resultados de este trabajo adecuada al público objetivo, como se detallará posteriormente. Asimismo, si bien se ha encontrado un servicio similar para filtrar por diversos criterios los municipios de Castilla y León \cite{repueblame}, o diversas iniciativas para descubrir pueblos de la España vaciada \cite{apadrina_un_pueblo}, \cite{pueblos_magicos}, no se ha encontrado ninguna alternativa para descubrir ciudades y pueblos españoles, sin importar su densidad de población, y cuya recomendación esté personalizada para cada usuario.
 \end{itemize}

\capitulo{2}{Objetivos del proyecto}

A continuación, se exponen los principales objetivos de este trabajo:

\begin{enumerate}
    \item \textbf{Construir un sistema de recomendación de municipios basado en las preferencias del usuario.} Como parte de este objetivo, se acotará la definición de \guillemotleft preferencias del usuario\guillemotright\space y la expectativa de cómo se han de manifestar y tratar. Además, queda establecido por el objetivo que habrá algún tipo de interacción que permitirá introducir datos y mostrar el resultado correspondiente.
    \item \textbf{Recopilar los datos suficientes como para crear un perfil adecuado de cada municipio.} Este objetivo abarca todas las fases de extracción, transformación y limpieza de datos, con el objetivo de garantizar una adecuada cantidad y calidad en los datos con los que se alimentará al sistema.
    \item \textbf{Crear un modelo que se comporte adecuadamente.} Para conseguir esta meta, el sistema deberá ser capaz de recomendar municipios similares y no necesariamente evidentes, cuya sugerencia realmente aporte una utilidad real, sea pertinente y coherente; es decir, se evaluará el rendimiento del sistema y se ajustará de manera que arroje las mejores métricas posibles.
    \item \textbf{Explorar los dos tipos de sistemas de recomendación fundamentales.} Estos son los basados en contenido y los basados en filtros colaborativos. Se procurará explorar ambas ideas con la intención de buscar los casos de uso que mejor se adaptan a cada una, viendo cuál es la que mejores resultados ofrece y analizando los retos encontrados y superados en ambas.
    \item \textbf{Ofrecer un sistema que pueda ser puesto en producción fácilmente.} El sentido de este objetivo es doble: Por una parte, garantizar la búsqueda de soluciones viables y eficientes, que permitan devolver el resultado de una recomendación fácilmente, sin una demora significativa que impida su uso de forma cómoda. Por otra parte, busca aplicar los conocimientos de las asignaturas de la rama de ciencia de datos, dado que será necesario almacenar grandes cantidades de información y operar con ellas en alguna de las soluciones de infraestructura escalable estudiadas.
    \item \textbf{Crear el sistema velando por las mejores prácticas que permitan su escalado y mejora continua.} Con este fin se pretende que todas las fases del proyecto –diseño, extracción, transformación, limpieza, modelado, programación, implementación y despliegue– apliquen las mejores prácticas posibles, cada una en su ámbito, para garantizar el escalado del sistema en el futuro –en el hipotético caso de que aumente el número de municipios, por ejemplo–, o su mejora continua –en el supuesto de que se desee hacer actualizaciones periódicas de los datos para mantener el modelo fiable y actual. Este objetivo abarca realizar una adecuada modularización del proyecto software, desacoplar en la mayor medida posible los datos de sus fuentes, utilizar estructuras de datos agnósticas a la presentación final, y aplicar los principios de responsabilidad única, por ejemplo, junto con emplear las mejores prácticas posibles en la implementación técnica de cada fase, como se describirá posteriormente.
    \item \textbf{Garantizar un adecuado nivel de seguridad jurídica y técnica en cuanto a la puesta en marcha se refiere.} Como parte integral del diseño y desarrollo del sistema, y aplicando los conocimientos de la parte del máster relacionada con las ramas de seguridad y derecho informáticos, se ha pretendido aplicar los principios de seguridad y privacidad por diseño y por defecto, incluidos en el Reglamento Europeo de Protección de Datos~\cite{aepd}, que buscan considerar desde el primer momento de ideación la información personal con la que tratará el sistema, intentando minimizarla en la medida de lo posible y restringirla a la imprescindible para funcionar adecuadamente. Esto es especialmente deseado, tal como aboga el Reglamento, en aquellas actividades que conlleven la creación de perfiles individuales o que procesen cantidades de datos de forma masiva, características que podría tener el sistema resultante de este trabajo.
    \item \textbf{Validar el sistema con usuarios potenciales reales.} Se buscará obtener retroalimentación de personas cercanas al autor que cumplan las condiciones para ser usuarios potenciales del resultado final del trabajo, a fin de validar cualitativamente diversos aspectos del mismo, como la utilidad, capacidad de interacción, tiempo de respuesta del sistema, facilidad de uso, satisfacción con la forma de organizar y presentar la información, coherencia de las respuestas ofrecidas frente a las esperadas, etc. Para ello, se hará uso de diversos conocimientos adquiridos en la asignatura \guillemotleft Visualización de Datos\guillemotright.
    
\end{enumerate}
\capitulo{3}{Conceptos teóricos}

A continuación, se exponen los principales conceptos teóricos cuya aparición en el presente trabajo es relevante.

\section{Sistema de recomendación}

Un sistema de recomendación, también llamado a veces recomendador automático –para reforzar la ausencia de intervención humana–, es una herramienta informática que opera como un sistema de filtrado de información, de forma que, para un conjunto de datos, devuelve los candidatos más prometedores con respecto a la función para la que fue diseñado. Es el caso de los mecanismos que recomiendan qué vídeo, serie o película ver a continuación, qué noticias relacionadas existen o qué artista musical nos puede gustar descubrir. La principal característica de estos sistemas es su capacidad para encontrar cosas que el usuario no está buscando activamente, porque ignora que existen, pero que encajan con su perfil y, por lo tanto, resultan útiles y agradables.

Tal como se vio en la asignatura \guillemotleft Técnicas de Aprendizaje Automático Escalables\guillemotright, se debe tener en cuenta factores como la relevancia (fundamental para ofrecer sugerencias útiles), la novedad o aleatoriedad donde corresponda (puede ser positivo priorizar elementos recientes, como noticias o canciones; o sorprender con resultados con un toque algo más aleatorio), y la diversidad (a fin de ofrecer resultados menos evidentes que los más lógicamente esperables).

Existen dos aproximaciones fundamentales: los sistemas de recomendación basados en contenido y los que están basados en filtros colaborativos. Recientemente, han aparecido soluciones híbridas que aúnan ambos esquemas para ofrecer resultados más robustos, donde los artículos se parecen individualmente pero las recomendaciones también cuentan con el respaldo de usuarios reales. En ambos casos se desea crear listas de los $k$ elementos cuya recomendación es más favorable.

\subsection{3.1.1. Sistema de recomendación basado en contenido}

En estos sistemas se analiza el contenido de los artículos a comparar (libros, canciones, series...), y se extraen sus características para poder utilizar medidas de similitud entre las propiedades encontradas que aporten una idea de la similitud entre los ítems que forman.

En ocasiones se construyen creando primero perfiles de los usuarios para agrupar los elementos más parecidos, con técnicas similares al \textit{clustering}, como variantes del algoritmo de k-vecinos más cercanos (\textit{k-nearest neighbors}).

La principal limitación de usar este enfoque es el llamado problema de arranque en frío (o \textit{cold start problem}, en inglés), que se caracteriza por no saber cómo actuar ante un caso nuevo, tanto de artículos como de usuarios. Por otra parte, si no queremos depender de los perfiles de los usuarios para encontrar elementos similares –lo que nos permite poder ofrecer recomendaciones incluso aunque no tengamos usuarios en el sistema o este aún no haya expresado ninguna preferencia histórica–, deberemos asegurarnos de que contamos con características suficientes (en número y en representatividad) como para poder encontrar parecidos razonables entre los artículos.

Aquí la definición de similitud no es agnóstica del contenido, sino totalmente dependiente; de hecho, las medidas de distancia empleadas variarán también en gran medida según el dominio del problema. Así, para vectores de números reales se usará habitualmente el coeficiente de correlación de Pearson –que mide la dependencia entre dos variables aleatorias continuas, sin necesidad de que los valores estén en la misma escala– o la similitud coseno –que mide la variación en el ángulo entre dos vectores–; mientras que para vectores binarios (asociados a características cualitativas) se usará la similitud de Jaccard –que mide el número de elementos comunes entre dos conjuntos–.

\begin{itemize}
    \item Coeficiente de correlación de Pearson. Siendo X e Y un par de variables aleatorias continuas, es el cociente entre la covarianza entre ambas y el producto de su desviación estándar:
\end{itemize}

$\rho_{X,Y}={\sigma_{XY} \over \sigma_X \sigma_Y} =\frac{\operatorname{Cov}(X,Y)}{\sqrt{\operatorname{Var}(X)\operatorname{Var}(Y)}}$

\begin{itemize}
    \item Similitud coseno. Siendo A y B dos vectores, la similitud coseno es su producto escalar dividido por el producto de sus normas:
\end{itemize}

${\displaystyle {\text{similitud}}=\cos(\theta )={\mathbf {A} \cdot \mathbf {B} \over \|\mathbf {A} \|\|\mathbf {B} \|}={\frac {\sum \limits _{i=1}^{n}{A_{i}B_{i}}}{{\sqrt {\sum \limits _{i=1}^{n}{A_{i}^{2}}}}{\sqrt {\sum \limits _{i=1}^{n}{B_{i}^{2}}}}}}}$

\begin{itemize}
    \item Índice de Jaccard. Siendo A y B dos conjuntos, este índice se calcula dividiendo el número de elementos que se encuentran en su intersección, entre el número de elementos que se encuentran en la unión:
\end{itemize}

${\displaystyle {\mathcal {J}}={\frac {\mid A\cap B\mid }{\mid A\cup B\mid }}}$

\subsection{3.1.2. Sistema de recomendación basado en un filtro colaborativo}

Otra manera de evitar las limitaciones de los sistemas de recomendación basados en contenido es utilizar la \textit{sabiduría del grupo} (o \textit{wisdom of the crowd} en inglés), que nos permite ofrecer sugerencias comparando los gustos implícitos (clicks, visitas...) o explícitos (valoraciones, opiniones...) de diferentes usuarios; de forma que, si dos usuarios se parecen en sus perfiles, podremos asumir que existirá una posibilidad significativa de que lo que le ha gustado anteriormente a uno, le guste al otro.

Además, este tipo de sistemas presenta una mayor independencia del dominio, dado que se basa en evaluaciones agnósticas de las características de los elementos que componen el conjunto; y una mayor capacidad de explicación, dado que las preferencias se asocian automáticamente, en lugar de depender de las características extraídas.

Entre los retos que presentan estos sistemas nos encontramos su escalabilidad y dispersión de la matriz de gustos, dado que un usuario valorará habitualmente muy pocos elementos del total de elementos existentes; la privacidad, ya que los usuarios no siempre querrán que se les asocie con los elementos con los que interactúen; y la dificultad para recomendar a usuarios poco comunes.

Los principales métodos de esta familia se basan en memoria y en factorización de matrices. Los primeros buscan similitudes entre usuarios o artículos (filas o columnas, respectivamente de la matriz de gustos) para después aplicar algún algoritmo de la familia de los k-vecinos más próximos, pero tienen complejidades temporales y espaciales elevadas cuando crecen los datos. Por otra parte, los segundos se basan en modelos para implementar algún algoritmo de factorización de matrices, como ALS (siglas de \textit{Alternating Least Squares}, Mínimos Cuadrados Alternos, en inglés), que permite reducir una matriz de rango alto a dos matrices de rango menor cuyo producto se aproxima a la original y donde el cálculo de similitud entre usuarios o artículos es inmediato y se resuelve con un producto escalar de vectores fila o columna, respectivamente. Para estimar el error global se puede minimizar la raíz del error cuadrático medio, o \textit{Root Mean Squared Error} (RMSE, por sus siglas en inglés), que consiste en la raíz de la media del cuadrado de la diferencia entre los valores pronosticados y los valores observados. Siendo $n$ el número de observaciones, $R__ij$ el valor observado, y $U_i * I_j^T$ el valor pronosticado:

RMSE = $\sqrt{\frac{1}{n}\Sigma_{i=1}^{n}{(R_ij - U_i * I_j^T)}^2}$

\section{\textit{One Hot Encoding}}

Es un tipo de codificación de datos ampliamente usado en tratamiento de características categóricas donde el conjunto de posibles valores se va a representar por un grupo de bits, de forma que cada valor posible se representará con el bit correspondiente en valor alto, y todos los demás en valor bajo.

Por ejemplo, si quisiéramos representar la presencia o ausencia de palabras en un texto de un corpus, crearíamos una matriz con tantas filas como documentos componen el corpus, y tantas columnas como valores posibles tiene nuestro diccionario, marcando con un valor binario la presencia o ausencia de cada palabra en cada texto.

\section{Discretización}

Es el proceso de convertir funciones continuas en homólogas discretas, habitualmente necesario al digitalizar señales analógicas o tratar con variables cualitativas. Puede conseguirse creando intervalos en el rango de valores que toma la función, y aproximando los valores al intervalo más cercano~\cite{discretizacion}.

Siempre que se discretizan valores continuos, existe lo que se conoce como error de discretización, que deberá procurar reducirse a una cantidad aceptable para el modelado correspondiente.

\section{Normalización}

Este concepto tiene amplios significados en el campo de la estadística, pero por lo que a este trabajo respecta, nos referiremos por \guillemotleft normalización\guillemotright\space a la conversión de escala de la distribución de una variable para poder realizar comparaciones adimensionales entre diversos elementos o con respecto a sus promedios u otras propiedades estadísticas. De esta manera, es posible comparar valores naturalmente distintos, procedentes de diferentes variables medidas con diversas unidades~\cite{normalizacion, normalizacion2}.

Existen diversos tipos de normalización según su caso de aplicación:

\begin{itemize}
    \item Puntuación estándar: $\frac{X - \mu}{\sigma}$
\end{itemize}

Es  la diferencia entre el valor observado y la media, dividido entre la desviación típica. Se usa para normalizar errores cuando se conocen los parámetros de población, que debe seguir una distribución normal.

\begin{itemize}
    \item T de Student: $\frac{X - \overline{X}}{s}$
\end{itemize}

Normaliza los residuos cuando se pueden estimar los parámetros de población pero no se conocen.

\begin{itemize}
    \item Coeficiente de variación: $\frac{\sigma}{\mu}$
\end{itemize}

Se aplica en distribuciones positivas como la exponencial o la de Poisson, y utiliza la media como medida de escala para normalizar la dispersión.

\begin{itemize}
    \item Característica \textit{scaling}: $X' = \frac{X - X_{min}}{X_{max}-X_{min}}$
\end{itemize}

También se conoce como normalización basada en la unidad, dado que convierte los valores a valores del intervalo [0, 1].

\section{Organización territorial española}

En España existen actualmente 8131 municipios~\cite{municipio}. Los municipios son, según la Ley reguladora de las Bases del Régimen Local, las entidades locales básicas de la organización territorial del Estado~\cite{ESP_wiki}. De acuerdo con esta ley, los municipios poseen personalidad jurídica y capacidad para el cumplimiento de sus fines, contando con los elementos de su territorio, población y organización. Los municipios se organizan en provincias, y estas en Comunidades Autónomas, a excepción de las dos ciudades autónomas de Ceuta y Melilla, que no disponen de ningún municipio asociado ni pertenecen a otra comunidad autónoma.

Por su parte, existen 61 578 entidades singulares de población, que el Instituto Nacional de Estadística define como \textit{«cualquier área habitable de un término municipal, habitada, o excepcionalmente deshabitada, claramente diferenciada dentro del mismo y que es conocida por una denominación específica que la identifica sin posibilidad de confusión»}. Las entidades singulares pueden estar formadas por varios núcleos de población, y/o tenerla diseminada, englobando edificaciones que no pueden considerarse núcleo. Además, cada entidad posee una calificación tradicionalmente reconocida, como villa, lugar, aldea, urbanización, poblado, caserío, monasterio...

Es importante señalar que la división del municipio en entidades singulares y colectivas (agrupación de varias entidades singulares, o de partes de ellas) de población no goza de carácter oficial, pero sí gran tradición. No es infrecuente debido al alto número de entidades existentes conocer a una o a varias, e incluso presuponer que se trata de un municipio en sí mismo. Esto es especialmente notable en la mitad norte peninsular, a causa de la dispersión geográfica y la orografía, que junto con razones históricas motiva la existencia de numerosos diseminados, muchos de ellos actualmente deshabitados~\cite{INE_entidad_singular}.

Desde el año 1981, el Instituto Nacional de Estadística etiqueta cada núcleo de población con un código numérico formado por 11 cifras que, comenzando por la izquierda, tiene la siguiente forma:

\begin{itemize}
    \item Las dos primeras cifras corresponden a la provincia.
    \item Las tres siguientes (tercera, cuarta y quinta) corresponden al municipio.
    \item La sexta y séptima corresponden a la entidad colectiva, si la hubiera. De no ser el caso, se usarán dos ceros.
    \item La octava y novena cifras codifican la entidad singular dentro de la entidad colectiva o municipal, según corresponda.
    \item La décima y undécima se refieren a los núcleos de población dentro de la entidad singular, usando \guillemotleft99 \guillemotright\space para diseminados.
\end{itemize}

Los códigos fueron actualizados en el año 1991 para mantener un orden alfabético, y corresponde a los ayuntamientos revisar, al menos anualmente, la relación de entidades y núcleos de población, para ser remitida al Instituto Nacional de Estadística, que las publica con la misma frecuencia. En caso de que se incorporen nuevos núcleos, se les asignará un código correlativo al último asignado, no reutilizando nunca los códigos de los núcleos que pudieran desaparecer. Conviene señalar que es posible que existan municipios o incluso entidades singulares de población con el mismo nombre –de hecho, ocurre en varios casos–, por lo que la identificación debe hacerse con el código del INE siempre que sea posible. Sin embargo, debido al alto número de instituciones que no emplean este código, la desambigüación debe hacerse con la provincia de pertenencia, donde no existen municipios duplicados.

\section{Datos abiertos}

La filosofía de datos abiertos (u \textit{Open Data}, en inglés) busca liberar determinados tipos de datos, especialmente concernientes a instituciones públicas y privadas, sin que suponga un menoscabo de las libertades individuales, de forma libre; esto es, sin limitaciones de derechos de autor, patentes, regalías u otro tipo de mecanismos de control, amparándose en el derecho a saber de la ciudadanía, la transparencia como parte de un buen gobierno público y corporativo, y en la utilidad pública de herramientas que surgen de hacer la información accesible y reusable.~\cite{datos_abiertos}

Esto último se consigue con una limpieza previa de los datos, donde el organismo titular de los mismos se encarga de garantizar que los conjuntos de datos elegidos son aptos para ser liberados, tanto por su calidad (limpieza, corrección...), como por su capacidad jurídica para poder hacerlo (los datos están agregados, anonimizados o no se corresponden con personas físicas identificadas o identificables.

Además, es fundamental que las capacidades de adquisición y procesamiento técnicas de los datos sean grandes, con el objetivo de no obstaculizar el propósito mismo para el que los datos abiertos son diseñados: promover el desarrollo de contenidos periodísticos, soluciones técnicas, productos, servicios, etc. gratuita o comercialmente, que beneficien a la sociedad basándose en la idea de aprovechar datos, habitualmente de titularidad pública o considerados como de interés público (como datos de consumo, datos de calidad de aire y aguas, datos de tráfico, datos de contagios, etc...), para realizar un análisis y extraer conclusiones valiosas.

Esto se materializa en servicios que idealmente cuentan con APIs (interfaces de programación de aplicaciones, por sus siglas en inglés) y/o capacidad de exportación de los datos en múltiples formatos habitualmente reconocidos fácilmente por máquinas, como los formatos CSV (valores separados por comas, por sus siglas en inglés), JSON (\textit{JavaScript Object Notation}, notación de Objetos JavaScript en inglés), XML (\textit{Extensible Markup Language}, lenguaje de marcado extensible en inglés) o XLSX (formato de hojas de cálculo de Microsoft Excel, el menos reusable de los comentados por tratarse de archivos binarios de software propietario comercial; pero no por ello menos usado en administraciones públicas).

\section{\textit{Web scraping}}

Traducido en español como \textit{raspado Web}, abarca el conjunto de medios técnicos que posibilitan la extracción automática de información de sitios Web. Habitualmente, estas técnicas se basan en programas informáticos que simulan la navegación de un usuario legítimo en Internet, realizando las llamadas a través del protocolo de transferencia de hipertexto y su versión segura (HTTP y HTTPS, por sus siglas en inglés), o imitando los movimientos de un humano a través de un navegador Web con interfaz gráfica.

Dado que se basan en imitar el comportamiento de un humano, o la llamada a un servidor de un navegador Web, los resultados no son siempre los deseados, ya que existen diversas medidas técnicas que los autores de un sitio Web pueden adoptar, bien para evitar esta técnica en sus versiones más factibles, o bien como parte de las decisiones del diseño del sistema implementado, que pueden provocar que falle el rastreo Web. Un ejemplo de esto último son las páginas basadas en \textit{scroll} (o deslizamiento vertical) infinito, o que completan en algún momento la carga de datos de manera dinámica tiempo después de haber terminado de enviar la página al cliente de red. Este enfoque complica la creación de un programa que extraiga la información, dado que es notablemente más factible comunicarse vía HTTP/S con un servidor, lo que normalmente provoca el cierre de la conexión cuando se recibe la respuesta, no siendo posible emular el desplazamiento vertical o interacción que realizaría un humano.

Por estos motivos, esta técnica se usa como último recurso para extraer datos, prevaleciendo siempre que sea posible el uso de APIs públicas de la institución, que están en sí mismas diseñadas para un uso mecanizado y automático, y donde los cambios suelen estar indicados mediante mecanismos de versionado.

Es importante resaltar que el raspado Web, por definición, consiste en extraer datos que no necesariamente fueron concebidos para su extracción, por lo que puede conllevar implicaciones legales. Actualmente, la legislación no es clara en este sentido, y la jurisprudencia ha dictaminado veredictos cuyo resultado depende en gran medida de las características del \textit{scrapper} y del uso que se fuera a dar a esos datos~\cite{scraping_legal, ryanair}.

Entre los usos más comunes del Web \textit{scraping} encontramos la monitorización de precios, la inyección de datos de otros sitios Web de forma adaptada a las necesidades del que los inyecta, la lectura de contenido desde fuentes externas o la recopilación de datos para enriquecer búsquedas, creando fragmentos accesorios conocidos como \textit{rich snippets}, que aportan información de diversas fuentes de confianza relacionadas con la búsqueda.

\section{\textit{Geocoding} y distancias}

La geocodificación, referida habitualmente en inglés como \textit{geocoding}, abarca el proceso de asignar coordenadas geográficas a puntos geográficos cualitativos (ciudades, pueblos, direcciones, edificios, lugares emblemáticos, puntos de interés, accidentes geográficos...). De esta manera, se transforma una dirección exacta o aproximada, en unas coordenadas que permiten su representación inequívoca en un plano de un Sistema de Información Geográfica~\cite{geocodificacion}.

Conociendo la posición exacta de dos puntos es posible calcular su distancia por diversos medios, siendo algunas de las más empleadas las siguientes:

\begin{itemize}
    \item Distancia por carretera. Es una de las más populares por su utilidad, dado que tiene en cuenta la disponibilidad de infraestructura y las características orográficas del terreno para poder comunicar dos puntos cualesquiera por vía terrestre.
    \item Distancia geodésica. Una línea geodésica es aquella de menor longitud entre dos puntos en una superficie dada, estando contenida en la propia superficie. Habitualmente nos referimos a esta distancia como la distancia \textit{en línea recta}, y no tiene en cuenta los accidentes geográficos, infraestructura disponible u otras características que limiten la comunicación ente los puntos.
\end{itemize}

Pese a que el cálculo de distancias entre dos puntos geolocalizados (identificados con coordenadas geográficas, habitualmente latitud y longitud) puede requerir de un tiempo de computación elevado –especialmente si se tienen en cuenta factores como la conveniencia de la ruta o posibles combinaciones con medios de transporte colectivo–, es un cálculo factible. Sin embargo, no ocurre lo mismo con determinadas distancias que los humanos consideramos aún más intuitivas, como la distancia de un punto a la costa o a la montaña.

Incluso aún siendo capaces de definir formalmente conceptos como \textit{costa} y \textit{montaña} –asumiendo, por ejemplo, lugares de altitud menor y mayor, respectivamente, a unos umbrales determinados sobre el nivel del mar– no es trivial elegir el punto costero o montañoso más cercano para poder calcular algún tipo de distancia geográfica.

Un posible enfoque podría ser calcular la distancia a todos los puntos que componen la frontera natural de un territorio con la costa, y elegir la menor, pero podría estar sujeto a problemas de solución arbitraria, como en el caso de playas interiores o costas incluidas en el territorio. En el escenario de sistemas montañosos, sería necesario contar con el perímetro de todos ellos, algo difícil de conseguir en el caso de las montañas de menor altitud o que se encuentran fuera de sistemas montañosos importantes. Además, en ambos supuestos, el coste lógico que supone realizar comparaciones para un gran número de pares de puntos es elevado.

Otro enfoque puede ser contar con datos geolocalizados de accidentes geográficos similares, como \guillemotleft playas\guillemotright\space o \guillemotleft picos montañosos significativos\guillemotright, y calcular la distancia hasta los mismos. De nuevo, se puede caer en el error de asimilar una \guillemotleft playa\guillemotright\space con una \guillemotleft costa\guillemotright, teniendo problemas similares en el caso de las montañas (dado que se forzaría a escoger solo los picos más significativos, que son los que cuentan con denominación y geolocalización conocidos), incurriendo en costes altos de igual manera y dependiendo notablemente de la cantidad y calidad de los datos auxiliares etiquetados.

Una de las formas de resolver problemas de distancias en tiempos razonables consiste en emplear heurísticas para aproximar la solución del problema. Por ejemplo, se podría tomar la heurística de considerar que los puntos pertenecientes a subconjuntos territoriales de ámbito superior que limitan con el hito geográfico de interés presentan la menor distancia posible, y que esta aumenta conforme aumenta la distancia de otras agrupaciones territoriales a las anteriores. De esta forma, los municipios pertenecientes a provincias costeras, serán los que menor distancia a la costa tendrán, los municipios de provincias colindantes a las costeras presentarán la segunda menor distancia a la costa, los municipios de provincias colindantes a las anteriores, la tercera menor distancia, y así sucesivamente, teniendo en cuenta que ninguna agrupación puede tener dos asignaciones de distancia distintas. En el caso de la distancia a la montaña, se podría tomar una heurística similar con ligeras modificaciones según el dominio del problema.

\section{\textit{Paas}: Plataforma como servicio}

Una plataforma como servicio, o \textit{PaaS} por sus siglas en inglés, es un conjunto de servicios y herramientas proporcionadas por un proveedor de computación en la nube (o \textit{cloud computing}) que permiten abstraer la infraestructura subyacente (controlada por el proveedor) y ofrecer a los desarrolladores de aplicaciones un entorno de despliegue de aplicaciones con diversas prestaciones que difieren de las capacidades de una instalación en sistemas propios (o que abaratan significativamente sus costes)~\cite{paas}.

Dado que los detalles de infraestructura son abstraídos para el programador, este puede despreocuparse de la escalabilidad, velocidad o instalación de actualizaciones de la plataforma, que será configurable, pero dependerá en última instancia del proveedor de la computación en la nube. Además, el uso de APIs le permiten comunicarse con el sistema como si realmente este fuera conocido para él, por lo que los resultados a nivel de usuario son equivalentes.

Ejemplos de plataformas como servicio son App Engine, de Google; Azure, de Microsoft; AWS Lambda, de Amazon; o Heroku, de Salesforce.
\capitulo{4}{Técnicas y herramientas}

A continuación, se presentan las principales técnicas y herramientas usadas en la construcción de este trabajo.

\section{Metodología de trabajo}

Como se ha estudiado en la asignatura \guillemotleft Técnicas de Aprendizaje Automático Escalables\guillemotright, existen diversas metodologías para proyectos de Big Data o ciencia de datos en grandes cantidades. En este punto, se procede a hacer un breve análisis de algunas de las más populares, y se indica también cuál ha sido la elegida para realizar este trabajo y por qué:

\subsection{4.1.1. KDD (\textit{Knowledge Discovery in Databases})}

KDD, del inglés \guillemotleft Extracción de conocimiento en bases de datos\guillemotright, se originó en la comunidad investigadora, y distingue entre etapas de identificación, preprocesado y modelado. Pese a que existe una retroalimentación entre etapas, pone el foco en las etapas de preprocesado, evaluación y presentación del conocimiento adquirido, en línea con el peso mencionado anteriormente que estos pasos tienen en cualquier proyecto de Big Data. Los procesos de esta metodología se pueden resumir en cinco etapas de la siguiente manera:

\begin{enumerate}
    \item Selección de datos a explorar. En primer lugar, elegiremos los datos que queremos integrar en la base de datos.
    \item Preprocesamiento de datos. Abarca las etapas de limpieza, donde se eliminará el ruido u otros datos incoherentes; la integración, que consiste en la unificación de datos procedentes de diversos orígenes; y la selección final de datos, que toma de la base de datos existente en este punto los datos relevantes (es decir, un subconjunto de variables) para la exploración del fenómeno a estudiar.
    \item Transformación de datos. Dentro de esta etapa, los datos adquieren formas apropiadas para su posterior procesado. Ejemplos de estas tareas son la creación de sumatorios, el uso de funciones de agregación, u otras similares para resumir los datos.
    \item Minería de datos. Consiste en buscar patrones de interés o representativos en los datos por medio de diversas técnicas.
    \item Evaluación de patrones e interpretación del conocimiento. Los últimos pasos de esta metodología, y que condicionan una posible regresión al comienzo, abarcan identificar los patrones que realmente aportan información novedosa o interesante, junto con su adecuada visualización para poder representar este nuevo conocimiento extraído al usuario.
\end{enumerate}

\subsection{4.1.2. SEMMA (\textit{Sample, Explore, Modify, Model and Access})}

La metodología SEMMA (muestrear, explorar, modificar, modelar y acceder, por sus siglas en inglés) se basa en muestrear la base de datos principal, asumiendo que un procesamiento completo es complejo y lento de llevar a cabo. Pone el foco en la transformación y modelado de los datos, omitiendo aspectos de negocio. Sus principales etapas son las siguientes:

\begin{enumerate}
    \item Muestrear. Consiste en tomar muestras de pequeño tamaño de la base de datos, con el objetivo de poder realizar manipulaciones de forma ágil.
    \item Explorar. Esta fase busca aumentar el entendimiento de los datos, refinando el proceso mientras se buscan anomalías, tendencias o patrones.
    \item Modificar. Abarca la creación, selección y transformación de variables para reducir su número, con la intención de contar con datos relevantes.
    \item Modelar. Busca crear el modelo que mejor satisfaga los objetivos propuestos al inicio del proyecto.
    \item Acceder. Evalúa la confianza y utilidad de los resultados extraídos, provocando una repetición del proceso en caso de que los objetivos no se hayan cumplido.
\end{enumerate}


\subsection{4.1.3. CRISP-DM (\textit{Cross-Industry Standard Process for Data Mining})}

CRISP-DM (proceso estándar entre industrias para el minado de datos, por sus siglas en inglés) es la metodología más usada actualmente~\cite{crisp}, y se diferencia de las anteriores por considerar el proceso de negocio dentro del propio ciclo de vida de los datos. Así, sus principales fases son:

\begin{enumerate}
    \item Entendimiento de negocio. Desde una perspectiva de negocio, se busca entender los objetivos y requisitos del proyecto, con la intención de convertirlos en una definición de un problema de minado de datos, para diseñar posteriormente un plan inicial para alcanzar dichos objetivos. Esta fase es de gran importancia en esta metodología, ya que se determinan los antecedentes, metas estratégicas del proyecto y criterios de éxito. Además, se realiza un inventario de los recursos, un análisis de costes y beneficios estimados, se determinan los objetivos finales y se produce el plan preliminar mencionado. 
    \item Entendimiento de datos. Esta fase comienza con una colección de datos inicial y actividades para familiarizarse con los datos, explorarlos, identificar problemas de calidad, descubrir las primeras observaciones o detectar conjuntos interesantes para formar hipótesis a partir de información oculta.
    \item Preparación de datos. En este punto, se realizan las transformaciones adecuadas para construir el conjunto de datos final que compondrá el modelo. Abarca transformaciones, limpieza, reducción de variables o cambios de formato, y su resultado es un entregable apto para ser introducido en la herramienta de modelado elegida.
    \item Modelado. Se seleccionan y aplican varias técnicas de modelado, cuyos parámetros se ajustan hasta encontrar los valores óptimos. Habitualmente será necesario volver a la etapa de preprocesado anterior para adecuar los datos a una nueva técnica de modelado, ya que con frecuencia se requieren tratamientos o formatos especiales para ciertas técnicas.
    \item Evaluación. Desde un punto de vista de negocio, se revisan los pasos ejecutados para construir el modelo junto con el grado de consecución de los objetivos inicialmente propuestos y la revisión de los criterios marcados para su evaluación. Se busca encontrar un objetivo clave de negocio que no haya sido suficientemente satisfecho, y al final de esta fase se decide si los resultados son aptos para continuar o no.
    \item Despliegue. Al igual que en el resto de metodologías mencionadas, CRISP-DM entiende que el desarrollo del proyecto no termina normalmente con la creación del proyecto. Por el contrario, suele ser necesario interpretar y presentar los conocimientos obtenidos de una manera adecuada según los requisitos de negocio, pudiendo significar la creación de informes o cuadros de mando, la elaboración de procesos de minado repetibles, el acceso e interacción con sistemas que contengan los resultados obtenidos, etc.
    
\end{enumerate}

Una de las principales diferencias entre esta metodología y las anteriores, es la posibilidad de volver hacia una etapa anterior en cualquier paso de la metodología CRISP-DM, promoviendo la mejora continua flexible y ágil, necesarias en muchos entornos empresariales.

Esta flexibilidad para mejorar los conjuntos de datos constantemente, junto con la gran popularidad de la que goza y el énfasis que realiza en los requisitos de negocio y en los criterios de evaluación y aceptación –cuya definición y consecución es fundamental en cualquier proyecto real, como el que se ha buscado realizar con este trabajo–, han propiciado su elección para el mismo.

Por último, para administrar y monitorizar el progreso del desarrollo de este trabajo, se ha utilizado el gestor de tareas Todoist~\cite{todoist}, y las herramientas conectadas al sistema de control de cambios Git, como tickets o \textit{issues}, de los repositorios en línea de GitHub~\cite{github}.

\section{Fuentes de datos}

Debido a la importancia de los procesos de extracción, transformación y limpieza para el desarrollo de este trabajo, se considera relevante incluir las principales fuentes de datos consultadas o explotadas para el mismo:

\begin{itemize}
    \item Altitud por municipio. Fuente: OpenElevation~\cite{openelevation}.
    \item Centros de Atención Primaria, Centros Hospitalarios y Centros de Atención Urgente Extrahospitalaria. Fuente: Ministerio de Sanidad~\cite{salud}.
    \item Cobertura de banda ancha por municipio. Fuente: Ministerio de Asuntos Económicos y Transformación Digital~\cite{cobertura}.
    \item Datos climáticos por municipio. Fuente: OpenWeather~\cite{openweather}.
    \item Distancias a las capitales de provincias. Fuente: Cálculos por carretera de OpenTripPlanner~\cite{otp}, y en línea recta a través de la distancia geodésica de GeoPy~\cite{GeoPy} cuando la primera no se encontró.
    \item Extracto de información del municipio e imágenes del municipio. Fuente: Wikipedia~\cite{wikipedia_api}.
    \item Geoposicionamiento por municipio. Fuente: PositionStack~\cite{positionstack}.
    \item Lugares más significativos por municipio. Fuente: GeoApify~\cite{geoapify}.
    \item Nomenclátor de Entidades Singulares del INE, con municipio de pertenencia asociado. Fuente: INE~\cite{municipios_ine} y Francisco Ruiz~\cite{entidades}.
    \item Nomenclátor Geográfico de Municipios y Entidades de Población. Fuente: Centro Nacional de Información Geográfica~\cite{nomenclator}.
    \item Número de colegios o institutos por municipio. Fuente: Ministerio de Educación~\cite{colegios}.
    \item Número de ofertas de empleo disponibles por provincia. Fuente: Infojobs~\cite{infojobs}.
    \item Número de universidades por municipio. Fuente: Ministerio de Educación~\cite{universidades}.
    \item Precios medios de viviendas en venta, y precios medios de viviendas en alquiler por provincia. Fuente: Fotocasa~\cite{fotocasa}.
    \item Precios y número de viviendas en venta, y precios y número de viviendas en alquiler por municipio. Fuente: Idealista~\cite{idealista}, y Fotocasa~\cite{fotocasa}, usada en los casos en los que la primera falló.
    \item Provincias de España y sus capitales. Fuente: Libretilla~\cite{provincias}.
    \item Relieve de España. Fuente: El Orden Mundial~\cite{relieve}.
    \item Rentas brutas medias por municipio de la Comunidad Foral de Navarra. Fuente: Hacienda Foral~\cite{renta_navarra}.
    \item Rentas brutas medias por municipio de País Vasco. Fuente: Eustat~\cite{renta_euskadi, PIB_pais_vasco}.
    \item Rentas brutas medias por municipio. Fuente: Agencia Tributaria para todos los territorios excepto País Vasco~\cite{renta}, y la Comunidad Foral de Navarra~\cite{renta_navarra}.
    \item Rugosidad del terreno en España. Fuente: Fundación BBVA~\cite{rugosidad}.
    \item Tasa de actividad, paro y empleo por provincia y sexo. Fuente: INE~\cite{INE_empleo}.
\end{itemize}

Como se detallará posteriormente, en diversos momentos de las fases de ETL y exploración de datos ha sido necesario utilizar diversas utilidades de terceros para procesar parte de los datos, verificar expresiones o comprobar la integridad de los mismos. Estas han sido las siguientes:

\begin{itemize}
    \item Postman, de Postman~\cite{postman}. Utilizada para realizar peticiones HTTP y HTTPS a los servicios Web ajenos (para conocer el formato de la respuesta devuelta), y propios (para ahorrar peticiones innecesarias durante la programación de los \textit{scripts}).
    \item Comparador de listas, de Molbiotools~\cite{molbiotools}. Se ha utilizado para comparar listas y verificar la corrección de los cambios.
    \item Contador de líneas duplicadas en listas, de Shailesh N. Humbad~\cite{somacon}. Empleada para verificar la integridad de los datos.
    \item Conversor de CSV a JSON, de Flatfile~\cite{csv2json}. Empleada para crear archivos JSON para la Web a partir de los archivos CSV manejados habitualmente.
    \item Coordinates Plotter, de Darrin J. Ward~\cite{mapmaker}. Usada para explorar los datos de geoposicionamiento.
    \item Detector de duplicados, de dCode~\cite{dcode}. Ha sido usada para buscar duplicados en listas, con información complementaria a la del resto de herramientas.
    \item Excel, de Microsoft~\cite{excel}, y Numbers, de Apple~\cite{numbers}. Utilizadas para abrir los archivos en formato CSV, XLS y XLSX, y verificar su integridad y corrección.
    \item Herramienta de explicación, construcción y pruebas de expresiones regulares, de gskinner~\cite{regexr}. Empleada para construir y comprobar la corrección de diversas expresiones regulares.

    \imagen{regexr}{Manejo de expresiones regulares para detectar y corregir inconsistencias.}
    
    \item Herramienta para ordenar una lista alfabéticamente, de Scott Smind~\cite{alfabetico}. Utilizada para ordenar puntos de la memoria y anexos alfabéticamente.
    \item Supresor de líneas duplicadas en listas, de autoría anónima~\cite{dedupelist}. Usada para eliminar duplicados en listas intermedias.
    \item Supresor de líneas vacías en listas, de CodeBeautify~\cite{codebeautify}. Usada para eliminar líneas en blanco tras supresiones manuales o automáticas.
    \item Validador de datos en formato JSON, de CircleCell~\cite{jsonlint}. Utilizado para comprobar la integridad y corrección de datos en formato JSON.
    \item Visualizador de datos en formato JSON, de CodeBeautify~\cite{jsonviewer}. Utilizada para visualizar gráficamente el árbol jerárquico de un archivo en formato JSON.
\end{itemize}

\section{Programación}

Como se detallará posteriormente, el proyecto abarca \textit{scripts} de programación escritos en Python y JavaScript que extraen, transforman, limpian y procesan información de diversas fuentes (tanto propias como ajenas), tomando los datos de archivos CSV (archivos separados por comas, por sus siglas en inglés, o puntos y comas, en el caso de este trabajo), JSON (archivos de notación JavaScript), XML (lenguaje de marcado extensible) o páginas de Internet públicas; y devolviendo archivos en formato CSV o JSON.

Entre las bibliotecas de terceros usadas, destacan las siguientes, ordenadas alfabéticamente y respetando la grafía de los autores:

\begin{itemize}
    \item beautifulsoup4. Desarrollada por Leonard Richardson bajo licencia MIT~\cite{beautifulsoup4}, permite extraer fácilmente información de páginas Web mediante mecanismos de interpretación y selección avanzada de elementos HTML y XML. Se ha utilizado en los archivos de extracción de datos de origen Web.
    \item Flask. Publicada por Armin Ronacher con licencia BSD~\cite{flask}, es un micro \textit{framework} Web, o micro marco de desarrollo Web, para aplicaciones en línea basadas en Python. A diferencia de otros marcos, Flask es ligero y ofrece una gran flexibilidad en cuanto a arquitectura, herramientas, estructura de ficheros o dependencias, lo que lo convierte en una opción escalable y que se adapta bien a numerosos tipos de proyectos. Constituye la base sobre la que se ejecuta la aplicación Web final.
    \item Jinja2. Desarrollada por el mismo autor de Flask y bajo la misma licencia~\cite{jinja2}, Jinja es un motor de plantillas ligero y extensible, que permite crear páginas HTML dinámicas en tiempo real, facilitando enormemente la programación de sitios Web de gran tamaño. Se ha utilizado para generar las páginas HTML de la aplicación Web, con algunas particularidades que se detallarán posteriormente.
    \item pandas. Programada por el equipo de desarrollo de Pandas (\guillemotleft The Pandas Development Team\guillemotright\space), y publicada bajo licencia BSD~\cite{pandas}, proporciona estructuras de datos eficientes para manejar grandes cantidades de información, series temporales y estadísticas varias. Se ha utilizado para acceder y procesar la información de los archivos CSV y JSON.
    \item python-dotenv. Desarrollada por Saurabh Kumar con licencia BSD~\cite{py_dotenv}, permite utilizar variables de entorno en Python de fácilmente, evitando escribir secretos como claves de API o credenciales en el código en texto plano que queda registrado en el sistema de control de versiones. De esta manera, los secretos se almacenan codificados mediante un par clave-valor en un archivo \texttt{.env}, de variables de entorno. Este archivo se incluye en el fichero \texttt{.gitignore} para evitar su adición al repositorio. Tras incluir y llamar a la función principal de esta biblioteca, los secretos se cargan en el entorno del sistema para acceder a ellos de forma segura, quedando localmente en la máquina de destino y evitando su propagación en texto plano más allá del citado archivo. Se ha empleado para no exponer las credenciales de las APIs usadas en la fase de extracción Web y en algunos servicios del producto final. Esta es una de las medidas de seguridad estudiadas en la asignatura \guillemotleft Fundamentos de Ciberseguridad\guillemotright, y presente en diversos documentos de seguridad estándar, como el Plan Director de Seguridad del INCIBE~\cite{INCIBE} o el proyecto OWASP para la seguridad de aplicaciones Web~\cite{OWASP}.
    \item phpDotEnv. Desarrollada por Vance Lucas con licencia BSD (cláusula 3)~\cite{phpdotenv}, permite la misma medida de seguridad anterior en el servidor donde se ejecuta el código PHP.
    \item re. Este módulo, desarrollado por la fundación de software Python (\guillemotleft Python Software Foundation\guillemotright\space) forma parte de la distribución oficial de Python y utiliza su licencia propia (PSFL, por sus siglas en inglés)~\cite{re}. Permite operar con expresiones regulares, y se ha utilizado para buscar patrones, como por ejemplo precios o metros cuadrados en páginas Web, o paréntesis en los nombres de los municipios.
    \item requests. Publicada por Kenneth Reitz bajo la licencia Apache 2.0~\cite{requests}, esta biblioteca ofrece implementar la capa de red para comunicaciones HTTP y HTTPS de forma sencilla. Se ha utilizado para realizar las llamadas en los archivos de extracción Web.
    \item scikit-learn. Desarrollada por Andreas Mueller con licencia BSD~\cite{scikit}, es una serie de módulos para el aprendizaje automático y el minado de datos. Se ha utilizado para crear el modelo y los recomendadores.
    \item surprise. Incluida posteriormente como módulo de scikit-learn~\cite{scikit_surprise}, es una biblioteca desarrollada por Nicolas Hug con licencia BSD para sistemas de recomendación~\cite{surprise}. Se ha utilizado para desarrollar el recomendador de filtro colaborativo.
    \item sys. Mencionado por su gran presencia en el código, el módulo del sistema de Python permite acceder a variables de entorno y funciones que interactúan con el intérprete~\cite{sys}. Se ha utilizado para pasar argumentos a los \textit{scripts} por línea de comandos, indicando valores como los archivos de entrada o salida u otros parámetros de configuración de los mismos.
    \item wikipedia. Publicada por Jonathan Goldsmith con licencia MIT~\cite{wikipedia_py}, ofrece la API pública de Wikipedia a través de Python. Se ha utilizado para extraer la información de Wikipedia de los distintos municipios.
\end{itemize}

Se ha aprovechado la existencia de un servidor personal para generar archivos de imitación (\textit{mock}) de diversas fuentes a las que se quería acceder durante el proceso de extracción de datos (tanto vía API como vía HTTP), con la intención de probar los guiones de extracción escritos en Python y JavaScript contra ellos durante el desarrollo. De esta manera, se han utilizado estos archivos de imitación durante la fase de depuración de errores, verificando el correcto funcionamiento de los archivos antes de realizar las llamadas a los servicios de terceros, y evitando incurrir en gastos de cuotas innecesarios en los servicios accedidos mediante API, o en demasiadas llamadas a los servicios de terceros durante el proceso de \textit{Web scraping}.

Además, se han creado archivos de HTML, CSS y JavaScript para la visualización e interacción con el resultado del proyecto, que constituye una aplicación Web desarrollada con el \textit{micro-framework} Flask. Para su desarrollo, al igual que el resto de archivos de programación del proyecto, se ha empleado el editor de código Sublime Text~\cite{sublime}, y los navegadores Web Safari~\cite{safari}, Chrome~\cite{chrome} y Firefox~\cite{firefox}.

\section{Despliegue}

Para desplegar el resultado final del proyecto software se ha optado por la plataforma como servicio (PaaS, por sus siglas en inglés) App Engine, de Google~\cite{appengine}. Esta solución abstrae la infraestructura necesaria para desplegar una aplicación Web, dejando en manos del desarrollador la configuración del entorno a nivel de aplicación, y permitiendo que este pueda alojarse en los servicios de computación en la nube de Google a bajo coste. App Engine es una de las alternativas estudiadas en este máster dentro de la asignatura \guillemotleft Infraestructura para el Big Data\guillemotright, y ha sido elegida por la experiencia previa del autor con ella y su conveniencia frente a otras alternativas. Sin embargo, como se abordará posteriormente, no se descarta una posible migración en el futuro a otra alternativa, dada la facilidad para hacerlo entre las soluciones de computación en la nube existentes. Por otra parte, el dominio Web del proyecto ha sido adquirido con la empresa Soluciones Corporativas IP~\cite{dondominio}.

\section{Otras herramientas}

Como herramienta para la creación de la memoria y anexos en \LaTeX \space se ha empleado el portal en línea Overleaf~\cite{overleaf} y \TeX{}maker~\cite{texmaker}. Para la presentación de diapositivas se ha usado Keynote~\cite{keynote}, mientras que los diagramas de ingeniería de software se han creado mediante Visual Paradigm~\cite{visual_paradigm}. Los iconos empleados a lo largo del trabajo, el producto software o cualquier otro artefacto relacionado con este trabajo, y que no sean de fuente propia, están extraídos de Icons8~\cite{icons8}.
\capitulo{5}{Aspectos relevantes del desarrollo del proyecto}

En este apartado se recogen los hechos más significativos del desarrollo del proyecto.

Como se ha comentado en el capítulo anterior, la metodología de trabajo elegida ha sido CRISP-DM, que se ha usado como hilo conductor del trabajo. A continuación, se procede a comentar lo más representativo de cada etapa, junto con el trabajo que se ha llevado a cabo en cada una:

\section{Entendimiento de negocio}

Como parte de esta etapa, se definieron los ocho objetivos detallados en la segunda sección de esta memoria, y se definieron los siguientes requisitos para poder alcanzar dichos objetivos:

\begin{enumerate}
    \item El sistema deberá recoger las preferencias del usuario, en el sentido de sus gustos personales, para poder mostrar municipios que las satisfagan.
    \item Estas preferencias se expresarán a través de un formulario de diversas preguntas, que permitan segmentar los municipios candidatos hasta encontrar los que reúnan las características deseadas.
    \item Se implementará un sistema de recomendación basado en contenido, para lo que se deberá recoger la entrada del usuario explícitamente. Adicionalmente, se podrá recoger la entrada implícitamente en otros casos de uso.
    \item Se recogerá también el gusto o no de los usuarios por el resultado ofrecido, de manera que se puedan tener datos para construir un sistema de recomendación basado en filtro colaborativo.
    \item Se explorará un sistema de recomendación basado en filtro colaborativo.
    \item Se recopilarán datos significativos basándose en otros trabajos relacionados y en los que el autor considere relevantes, que permitan crear perfiles detallados para cada municipio, de entre los conjuntos y extracciones de datos disponibles para el ámbito del proyecto. En particular, siempre que sea posible se incluirán datos de: nombre completo del municipio, provincia y comunidad autónoma a las que pertenece, número de habitantes, superficie, densidad de población, disponibilidad de servicios educativos (colegios y universidades), disponibilidad de servicios sanitarios (centros de salud y hospitales), datos climáticos históricos básicos, datos de renta y empleo básicos, datos de venta y alquiler de viviendas en el municipio y en su provincia, cercanía a la capital de provincia, cercanía a la costa y a la montaña, altitud, datos de cobertura de Internet en el municipio o en su provincia, principales servicios de los que dispone, ubicación y extractos de información textual y fotográfica.
    \item El sistema será capaz de arrojar recomendaciones para un 70 \% de los casos probados, y su rendimiento se medirá por métricas ajustadas a la naturaleza de su modelo, buscando alcanzar una fiabilidad adecuada.
    \item Se diseñará el sistema de forma escalable; es decir, se considerará desde su diseño en posibles repeticiones automáticas del proceso de extracción, transformación y limpieza de datos, posibles adiciones o supresiones de la base de datos de municipios españoles, posibles incorporaciones, modificaciones o supresiones de datos para uno o varios municipios, posibles traducciones o posibles casos de uso similares a los existentes y potencialmente deseables.
    \item Se asignará responsabilidad única a los componentes informáticos que formen el trabajo.
    \item Se presentará un resumen de los datos de cada municipio al usuario.
    \item Se creará un sistema versionado, de fácil despliegue, actualización e interacción.
    \item Se creará un sistema con la privacidad y la seguridad en mente, que no expondrá secretos ni credenciales, que aplicará las mejores prácticas estándar para el ámbito del proyecto, que no recogerá datos personales y que contará con sistemas de contacto para comunicarse con el responsable del sistema.
    \item Se realizará una evaluación de la utilidad, usabilidad y rendimiento del sistema por parte de usuarios potenciales reales, a fin de detectar puntos de mejora y evaluar las decisiones de diseño e implementación adoptadas.
\end{enumerate}

Como criterios de éxito, además del valor de capacidad para obtener resultados establecido anteriormente, se requerirá una evaluación positiva, de al menos un 70 \% de satisfacción y un 50 \% de expectativa de uso por parte de los usuarios potenciales con los que se probará el sistema.

Para desarrollar el sistema se contará con dos ordenadores a disposición del alumno, el servidor personal mencionado anteriormente –con pila de tecnologías LAMP: Linux, Apache, MySQL y PHP–, y la infraestructura propia de Google –con entorno preparado para ejecutar aplicaciones Web de Python– donde se desplegará el resultado final del proyecto.

El desarrollo del trabajo tiene una duración estimada de alrededor de dos meses, y para completarlo se realizará un análisis de viabilidad previa, recolección y adecuación de los datos necesarios, integración de los datos y composición del modelo, evaluación del modelo e interpretación, presentación y despliegue del resultado final.

\section{Entendimiento de datos}

Una vez completado el análisis de viabilidad con éxito, gracias al descubrimiento de trabajos relativamente relacionados y de los conocimientos del autor en aplicaciones Web y su compatibilidad con otros sistemas basados en Python, se procedió a llevar a cabo la recolección y exploración inicial de datos. Para explicarla, se procede a detallar la función y contenido de cada componente software usado en el proceso de extracción. Nótese que, por fines de compartición de conocimiento y adecuación a los estándares del mundo académico y profesional, los nombres de archivos, variables y otros literales ajenos al modelo se expresan en inglés; mientras que los nombres de las columnas o variables propias de las fuentes de datos mantienen su nombre original en español. Por otra parte, en el caso de nombres escritos en idiomas co-oficiales se ha procurado mantener la convención del Instituto Nacional de Estadística, o la usada mayoritariamente entre las fuentes, en caso de conflicto entre varios nombres alternativos procedentes de diversas fuentes. Es importante recordar que se parte de la base de datos de municipios del INE en formato XLSX, que ha sido convertida a CSV y cuenta con las siguientes columnas: código numérico del municipio, nombre del municipio y provincia a la que pertenece. 

A continuación, se resumen los aspectos más relevantes de esta fase:

\begin{enumerate}
    \item Se hace uso de archivos no añadidos al repositorio público para contener las variables de entorno con secretos y credenciales de uso en los servicios de terceros usados.
    \item El cálculo de distancias a la costa y a la montaña desde cada municipio se ha realizado mediante heurísticas basadas en la presencia de costa en cada provincia, y en la importancia de sus sistemas montañosos, conforme a lo explicado anteriormente.
    \item Para los municipios cuya renta anual no es pública, bien por estar agrupada y promediada debido a su tamaño (menos de 1000 habitantes), o por secreto estadístico, se han tomado los valores oficiales más cercanos; esto es, el valor publicado de la agrupación para dichos municipios en esa provincia, y el valor publicado para el municipio más cercano, respectivamente.
    \item Se ha trabajado con el código del municipio asignado por el INE siempre que ha sido posible; sin embargo, debido a su baja adopción, ha sido necesario trabajar con dos claves primarias: municipio y provincia. Posteriormente, los datos de cada fuente que no lo tuvieran se han fusionado con el código del INE para poder realizar búsquedas unívocas y homogéneas.
    \item Dado que la mayoría de proveedores, de titularidad pública y privada, trabajan con el nombre del municipio y su provincia, es necesario tener en cuenta la frecuente normalización del nombre del municipio; es decir, la omisión de determinantes o su puesta al final del nombre, así como el predominio de la variante del idioma co-oficial local en la mayoría de los casos.
    \item Se emplean los datos de diversos terceros para obtener, por ejemplo, la geolocalización de cada municipio, su elevación, sus principales características o su tiempo atmosférico promedio en febrero y en julio. En particular, y como se explica detalladamente en el Anexo \ref{programador} \guillemotleft Documentación técnica de programación\guillemotright, en el caso de la localización de los municipios, se detectaron algunas inconsistencias en la fase de exploración de los datos; dado que el servicio localizó municipios homónimos fuera de España. Tras incluir el país como restricción para el servicio, la exploración de datos de este conjunto fue según lo esperado.
    \item Se ha realizado extracción de datos Web, o \textit{Web scraping}, para diversos conjuntos de datos, como los centros de educación superior o los datos inmobiliarios por municipio. Como se explica en el mencionado Anexo, esta extracción automatizada presentó diversos retos; como la necesidad de esperar un pequeño tiempo (uno o dos segundos) entre cada petición para no sobrecargar el servidor de destino, o la combinación de varias páginas o proveedores para obtener el resultado final.
    \item Se han creado diversos ficheros con funciones auxiliares, como fusionar conjuntos de datos hasta crear el modelo, extraer los nombres co-oficiales locales de los municipios, o eliminar columnas de diversos subconjuntos.
    \item El cálculo de la distancia por carretera entre cada municipio y la capital de su provincia se ha realizado mediante la herramienta OpenTripPlanner (OTP)~\cite{otp}. Básicamente, suministrándole el mapa de España, esta herramienta de software libre crea un servidor al que es posible realizar consultas de creación de rutas mediante una interfaz gráfica Web y una API HTTP. Este servidor se desplegó en un ordenador de más memoria RAM que el que se ha usado para desarrollar el trabajo, ya que el mapa requería al menos 28 GigaBytes, pero en la misma red doméstica que el principal del alumno. De esta manera, los archivos que realizan la extracción se conectan a través de la API HTTP al servicio, obteniendo los kilómetros que separan cada localidad de su capital de provincia.
    \item Se ha obtenido una lista de los veinte lugares más significativos de cada municipio, a través de la API del proveedor GeoApify~\cite{geoapify}, que presentaba la tarificación más atractiva para este propósito. Una vez extraídos, se realizan cálculos para ver qué categorías de las siguientes son las más repetidas: comercios, turismo, alojamientos, ocio, naturaleza, servicios, actividades, entretenimiento, hostelería y deporte, además de otras devueltas adicionalmente por el servicio. Posteriormente, se incorporan estos datos al modelo para tener información cualitativa de cada municipio.
    \item Los datos climáticos recogidos a través de la API de OpenWeather~\cite{openweather} son los siguientes:

    \begin{itemize}
        \item Promedio de las temperaturas mínimas en febrero.
        \item Promedio de las temperaturas máximas en febrero.
        \item Promedio de las temperaturas en febrero.
        \item Promedio de humedad en febrero.
        \item Promedio de la velocidad media del viento en febrero.
        \item Promedio de las precipitaciones mínimas en febrero.
        \item Promedio de las precipitaciones máximas en febrero.
        \item Promedio de las precipitaciones en febrero.
        \item Promedio de porcentaje de nubosidad en febrero.
        \item Promedio de número de horas de sol en febrero.
    
        \item Promedio de las temperaturas mínimas en julio.
        \item Promedio de las temperaturas máximas en julio.
        \item Promedio de las temperaturas en julio.
        \item Promedio de humedad en julio.
        \item Promedio de la velocidad media del viento en juliio.
        \item Promedio de las precipitaciones mínimas en julio.
        \item Promedio de las precipitaciones máximas en julio.
        \item Promedio de las precipitaciones en julio.
        \item Promedio de porcentaje de nubosidad en julio.
        \item Promedio de número de horas de sol en julio.
    \end{itemize}

    \item A través de la API de Wikipedia se ha extraído una breve descripción del municipio, así como algunas de las principales imágenes de su entrada en dicha enciclopedia. Estos textos e imágenes se usarán en el resultado final con fines ilustrativos, para que el usuario pueda conocer mejor el municipio recomendado y pueda aprender más sobre él. Con el fin de obtener imágenes de la localidad y su entorno, se excluyen arbitrariamente imágenes cuyo nombre involucre mapas, escudos, banderas, planos, etc.
    
\end{enumerate}

Una descripción exhaustiva de cada archivo involucrado en esta fase, y una explicación más detallada de los aspectos anteriores pueden encontrarse en el Anexo \ref{programador} \guillemotleft Documentación técnica de programación\guillemotright. 

Todos los datos han sido recogidos en diciembre de 2022, y hacen referencia a los datos más recientes disponibles a esa fecha, siendo las fechas más distintas a la fecha de recogida las de los datos:

\begin{itemize}
    \item Rentas brutas medias por municipio, pertenecientes al año 2019.
    \item Cobertura de banda ancha por municipio, perteneciente al año 2020.
\end{itemize}

El perfil estadístico de los datos que compondrán el modelo puede encontrarse en el Anexo \ref{diseno} \guillemotleft Especificación de diseño\guillemotright.

\section{Preparación de datos}

Una vez la exploración de los datos cuantitativos, que alimentarán posteriormente el modelo, fue exitosa en todos los conjuntos de datos, dio comienzo la fase de preparación de los datos.

En esta fase se eliminaron los marcadores de datos nulos o no encontrados (establecidos arbitrariamente como \guillemotleft -\guillemotright), y se sustituyeron por campos vacíos.

\imagen{modelo1}{Datos geográficos y de lugares de interés de los municipios antes de su transformación.}

Además, se convirtió la expresión de los lugares más frecuentes para un municipio de cualitativa, para la exploración de datos, a cuantitativa siguiendo la codificación \textit{one hot encoding}.

\imagen{modelo2}{Datos geográficos y de lugares de interés de los municipios después de su transformación.}

Antes de continuar, se creó una versión del modelo hasta el momento en formato JSON, mucho más conveniente a la hora de recuperar los datos de un municipio en particular para presentarlos al usuario, tanto por razones de comodidad como de rendimiento. Así, el sistema de recomendación devolverá el código INE del municipio, y sus datos se recuperarán de este archivo que, por este motivo, estará formado por un diccionario indexado por estos códigos.

Posteriormente, se eliminaron las variables cualitativas (nombre codificado del municipio, nombre tradicional del municipio y comunidad autónoma), y las coordenadas (campos de latitud y longitud) de cada municipio, y se normalizaron los datos para disponer de ellos de forma agnóstica a las unidades empleadas en cada dimensión. La normalización se hizo basándose en diversas variables, como la renta bruta media o la tasa de empleo; sin embargo, la que mejores resultados arrojó fue el tamaño de la población. Así, se utilizó la regresión de Ridge –que reduce los coeficientes introduciendo un término de penalización igual a la suma de los coeficientes cuadrados por un coeficiente de penalización– para normalizar las variables sobre la clase, y se multiplicaron los coeficientes devueltos por la estructura de datos.

Para todo ello se usan los archivos \texttt{normalize.py} y \texttt{create-model.py}.

\section{Modelado}

A la hora de hablar del modelo del trabajo, es necesario distinguir entre los dos enfoques que se han explorado:

\subsubsection{Sistema de recomendación basado en contenido}

Una vez se dispone del modelo normalizado, este puede procesarse para producir el sistema de recomendación basado en contenido. Aquí se ha optado por basarlo en el cálculo de la similitud coseno entre el municipio que constituirá la entrada y los restantes, con el objetivo de devolver el más cercano del espacio vectorial que forman los municipios representados por vectores según sus características.

Tras calcular la similitud coseno se ordenarán los valores de forma descendente, para seleccionar a continuación el resultado de mayor similitud, que será el municipio presentado al usuario. Para ello, el código del municipio devuelto será usado para acceder a la base de datos de los municipios en formato JSON.

\imagen{cosine}{Explicación de cómo la similitud coseno calcula la distancia entre vectores, devolviendo valores más altos cuanto más similares son; y más bajos cuanto más opuestos son~\cite{US3}.}

Así, se definió una función que toma un código INE como entrada y tras leer el CSV del modelo calcula la similitud coseno entre ese municipio y el resto, devolviendo el código del INE correspondiente al primer resultado. En la fase de despliegue, se preparó el sistema para recuperar la información del municipio a partir de su código INE y devolverla al usuario.

\subsubsection{Sistema de recomendación de filtro colaborativo basado en usuarios}

Además del sistema de recomendación basado en contenido, el autor deseaba construir un sistema de recomendación basado en un filtro colaborativo que se alimentara de las valoraciones (positivas o negativas de forma binaria, \guillemotleft me gusta / no me gusta\guillemotright) de los usuarios al recibir la recomendación de un municipio. La idea detrás de estos sistemas es recomendar elementos que han gustado a usuarios similares. Debido a la baja tasa de valoración habitual, suelen trabajar con matrices dispersas y necesitar grandes cantidades de valoraciones de usuarios para poder realizar recomendaciones adecuadas.

Dado que el dominio del problema no dispone de conjuntos de datos con valoraciones de usuarios y, el propio sistema, aunque diseñado para recogerlas, no ha podido coleccionar las suficientes durante su creación por los limitados recursos de un trabajo académico como este, el autor consideró crear su propio conjunto de datos ficticio imitando usuarios reales.

Para ello, creó 10 usuarios ficticios y valoró con cada uno a 100 municipios de forma positiva y a 100 municipios de forma negativa. Estos municipios estuvieron elegidos aleatoriamente aplicando los siguientes criterios para cada usuario, con la intención de crear perfiles de usuarios intuitivamente naturales:

\begin{itemize}
    \item El primer usuario muestra su valoración positiva por 100 municipios costeros aleatorios, y su valoración negativa por 100 municipios muy alejados de la costa aleatorios; es decir, representa a un usuario que prefiere los municipios costeros.

    \item El segundo usuario muestra su valoración positiva por 100 municipios muy alejados de la costa aleatorios, y su valoración negativa por 100 municipios costeros aleatorios; es decir, representa a un usuario que prefiere los municipios alejados de la costa.

    \item El tercer usuario muestra su valoración positiva por 100 municipios cercanos a la montaña aleatorios, y su valoración negativa por 100 municipios muy alejados de la montaña aleatorios; es decir, representa a un usuario que prefiere la montaña.
    
    \item El cuarto usuario muestra su valoración positiva por 100 municipios muy alejados de la montaña aleatorios, y su valoración negativa por 100 municipios de montaña aleatorios; es decir, representa a un usuario que prefiere los municipios alejados de la montaña.

    \item El quinto usuario muestra su valoración positiva por 100 municipios de menos de 10 000 habitantes aleatorios, y su valoración negativa por 100 municipios de más de 10 000 habitantes aleatorios; es decir, representa a un usuario que prefiere los pueblos a las ciudades.

    \item El sexto usuario muestra su valoración positiva por 100 municipios de más de 100 000 habitantes aleatorios, y su valoración negativa por 100 municipios de menos de 10 000 habitantes aleatorios; es decir, representa a un usuario que prefiere las ciudades a los pueblos.

    \item El séptimo usuario muestra su valoración positiva por 100 municipios con lugares naturales aleatorios, y su valoración negativa por 100 municipios sin lugares naturales aleatorios; es decir, representa a un usuario que prefiere la naturaleza.

    \item El octavo usuario muestra su valoración positiva por 100 municipios con lugares de comercios aleatorios, y su valoración negativa por 100 municipios sin lugares de comercio aleatorios; es decir, representa a un usuario que prefiere las zonas comerciales.

    \item El noveno usuario muestra su valoración positiva por 100 municipios con menos de un 50 \% de conectividad de fibra óptica a más de 100 megabytes aleatorios, y su valoración negativa por 100 municipios con más de un 50 \% de dicha conectividad aleatorios; es decir, representa a un usuario que prefiere zonas poco conectadas a Internet.

    \item El décimo usuario muestra su valoración positiva por 100 municipios con más de un 50 \% de conectividad de fibra óptica a más de 100 megabytes aleatorios, y su valoración negativa por 100 municipios con menos de un 50 \% de dicha conectividad aleatorios; es decir, representa a un usuario que prefiere zonas conectadas a Internet.

\end{itemize}

Este proceso se realizó a través del \textit{script} \texttt{create\_ratings.py}. A continuación, una vez creados los usuarios anteriores y las 2000 valoraciones iniciales, se usó la librería \texttt{surprise}, de scikit-learn, para tratar con ellos. Tras especificarle la escala de las opiniones (valores enteros entre el 0 y el 1), y especificar las columnas del conjunto (identificador del elemento, identificador del usuario y valoración), se usó primeramente KNNWithMeans~\cite{scikit_surprise} como algoritmo de filtrado colaborativo básico, que tiene en cuenta las valoraciones promedio de cada usuario.

Tras separar el conjunto en las particiones de entrenamiento (80 \% del conjunto) y prueba (20 \% restante), y entrenar el modelo y predecir con este algoritmo, se usó ALS (mínimos cuadrados alternos, por sus siglas en inglés) para después evaluar la precisión a través del RMSE (error cuadrático medio, por sus siglas en inglés), como se describirá en la fase de evaluación.

A partir de aquí, el sistema funciona de forma análoga al sistema basado en contenido, devolviendo identificadores de municipios recomendados para uno dado.

Todo ello se realizó en el archivo \texttt{users\_recommendations.py}.

Sin embargo, y como se abordará en el apartado de evaluación, este enfoque no dio los resultados deseados, teniendo un rendimiento inferior al esperado –quizás debido a la generación artificial de los usuarios–. Por este motivo, se decidió explorar el segundo tipo de sistemas de recomendación de filtro colaborativo.

\subsubsection{Sistema de recomendación de filtro colaborativo basado en ítems}

En estos sistemas, las recomendaciones se hacen teniendo en cuenta la similitud de los elementos valorados por otros usuarios con respecto a las valoraciones del usuario interesado. Se buscan elementos similares a los gustos comunes de usuarios cercanos. Si bien este método puede funcionar peor para usuarios con gustos muy específicos (usuarios de nichos, segmentos caracterizados por pocos usuarios con gustos comunes), el filtro colaborativo basado en ítems habitualmente ofrece recomendaciones mejores y más rápidas~\cite{filtro_colaborativo_2}. En el caso de este trabajo, se consideró que los municipios españoles reúnen las suficientes características como para poder ser divididos en un número relativamente pequeño de grupos, lo que favorece la ausencia de usuarios de nicho, por lo que se consideró un enfoque apropiado.

Así, para predecir la valoración (uno en caso de que la valoración sea positiva; es decir, el usuario expresa que le gusta el municipio, y cero en caso de que este exprese su rechazo) de un lugar para un determinado usuario, calcularemos la predicción de la valoración de un municipio $i$ por un usuario $u$ de la siguiente manera:

$$pred(u, i) = \frac{\sum_{v \in usuarios \ que \ han \ valorado \ i} sim(u, v) \cdot r(v, i)}{\sum_{v \in usuarios \ que \ han \ valorado \ i} |sim(u, v)|}$$

Donde:

$pred(u, i)$: es la predicción de la puntuación del usuario $u$ para el municipio $i$.

$sim(u, v)$: es la similitud entre los usuarios $u$ y $v$, medida por su similitud coseno.

$r(v, i)$: es la evaluación del municipio $i$ por el usuario $v$, que toma 1 si le ha gustado, y 0 en caso contrario.

Este procedimiento es análogo a la frase \guillemotleft la gente a la que le gusta el elemento X, como el usuario actual, también tiende a valorar positivamente el elemento Y, que el usuario aún no ha valorado, por lo que debería probarlo\guillemotright. El filtrado colaborativo basado en ítems tiene menos errores que el filtrado colaborativo basado en usuarios y, como su modelo menos dinámico se calcula con menos frecuencia, se puede representar en una matriz de menor tamaño, por lo que su rendimiento es superior al de los sistemas de filtrado colaborativo basados en usuarios~\cite{filtro_colaborativo_3}.

Además, para ofrecer la mayor cantidad de resultados posible y la mejor experiencia de usuario, al poder descubrir una elevada cantidad de municipios, se decidió hacer uso de una de las técnicas más usadas en recomendadores por su elevado éxito: Los sistemas de recomendación híbridos. De esta manera, a la lista de municipios sugerida según el cálculo anterior –es decir, según los gustos de otros usuarios similares–, se añadieron municipios similares por contenido a los ya valorados positivamente por el usuario. Como se detallará en la sección de despliegue, esta lista se computa una vez al día para todos los usuarios que alguna vez han valorado algún artículo en el sistema.

\section{Evaluación}

De nuevo, debemos hacer distinción entre los dos sistemas usados:

\subsubsection{Sistema de recomendación basado en contenido}

Se usaron diversas métricas para comprobar el funcionamiento de este sistema:

En primer lugar, la puntuación de Ridge devuelve el coeficiente de determinación de la predicción, siendo mejor cuanto mayor es, y estando su límite superior en 1. Se basa en la creación de un modelo de regresión lineal mediante cuadrados mínimos con norma L2. Su valor no está acotado inferiormente y puede tomar valores negativos, ya que el modelo puede ser arbitrariamente malo. Entre esta fase y su habitual vuelta a la fase de modelado anterior, se probaron las siguientes clases como base del modelo: renta bruta media, que obtuvo una puntuación de 0,62; tasa de actividad, que obtuvo una puntuación de 0,85; tasa de empleo, que obtuvo una puntuación de 0,91; y población, que obtuvo una puntuación de 0,96, siendo por tanto la variable elegida.

Por otra parte, se tomó un subconjunto de 1000 municipios, que se dividieron tomando el 80 \% de ellos en un conjunto de entrenamiento, y un 20 \% en un conjunto de prueba. Tras realizar un etiquetado manual de municipios parecidos, se usó el conjunto de prueba para calcular su RMSE, obteniendo un valor de 0,19, dentro de los objetivos del trabajo.

Además, se realizaron diversas pruebas para comprobar la fiabilidad del sistema de forma cualitativa, y ver si los resultados devueltos eran coherentes y se adecuaban a la expectativa del sistema: Se observó cómo, en todos los casos (100 pruebas formales), el sistema recomendó municipios de características similares; por ejemplo, municipios medianos - grandes como recomendación sugerida para otros municipios del mismo tamaño, municipios pequeños como recomendación sugerida para otros municipios del mismo tamaño, municipios costeros como recomendación sugerida para otros municipios costeros y municipios cercanos a la montaña para otros municipios cercanos a la montaña.

En general, la calidad de las recomendaciones observadas fue muy satisfactoria, tanto para el autor como para los sujetos que evaluaron cualitativamente el sistema una vez desplegado, como se describe posteriormente.

\subsubsection{Sistema de recomendación de filtro colaborativo basado en usuarios}

Por otro lado, el sistema de recomendación de filtro colaborativo basado en usuarios no alcanzó una fiabilidad adecuada en sus predicciones. La estimación del error cuadrático medio tras aplicar el algoritmo KNNWithMeans basado en la similitud coseno era demasiado grande (0,44), por lo que se decidió volver a la fase anterior y usar como método BaselineOnly con ALS y 3 ciclos completos (\textit{epochs}). De esta manera, se obtuvo un RMSE mucho menor, de 0,04, por lo que en principio los resultados eran más prometedores.

Sin embargo, al realizar las predicciones del conjunto de evaluación, los resultados no fueron satisfactorios. Las probabilidades de recomendar o no un municipio para un usuario similar a otro no superaban el 52 \%, quedando demasiado cerca a criterio del autor de una recomendación puramente azarosa.

Se pensó que podría deberse a un tamaño demasiado pequeño en el conjunto de muestra inicial (2000 valoraciones), por lo que se optó por aumentar el número de valoraciones de cada usuario hasta las 1000 por tipo de valoración, resultando en 20 000 valoraciones. Pese a ello, los resultados tras entrenar, probar y evaluar el modelo fueron similares, no alcanzando un 53 \% de precisión en las recomendaciones. Tras aumentar el conjunto de datos de nuevo por diez, creando 10 usuarios nuevos réplica de los anteriores y contar con 200 000 valoraciones, solo se aumentó el tiempo de cálculo necesario, obteniendo los mismos resultados en cuanto a su precisión, muy similar a la de recomendaciones aleatorias.

Se consideró que este hecho puede haber estado relacionado con el tamaño de la muestra (pocos usuarios y pocas valoraciones) o con los criterios arbitrarios elegidos para separar a los usuarios, que no reflejen adecuadamente el comportamiento y valoraciones de usuarios reales.
 
\subsubsection{Sistema de recomendación de filtro colaborativo basado en ítems}

Por su parte, el sistema de filtrado colaborativo basado en artículos obtuvo un RMSE de 0,25 en el conjunto de prueba, generado con 100 valoraciones aleatorias sobre un conjunto de 300 municipios aleatorios, parecidos por pares, siguiendo una proporción 80-20 para cada conjunto de datos (entrenamiento y prueba, respectivamente).

Para asegurar una calidad suficiente, las recomendaciones se tienen en cuenta si el valor de la predicción supera el umbral de 0,7. Además, al igual que en el recomendador de contenido, se realizó la misma evaluación cualitativa formal, consistiendo en 100 pruebas a 10 usuarios distintos (10 pruebas con cada usuario).

Tras alimentar el sistema con sus valoraciones, debían juzgar si las predicciones les resultaban útiles o no. Todos ellos consideraron que las predicciones mostraban municipios similares a aquellos por los que habían expresado su preferencia en el pasado, considerando la prueba satisfactoria.

\section{Despliegue}

Siguiendo con los objetivos de este trabajo, se encuentra una forma de interacción cómoda, centrada en el usuario, que le permita descubrir municipios en base a sus gustos y a los de otros usuarios. Para ello, se considera crear un sitio Web público como la mejor forma de lograrlo. El sitio Web está disponible en \url{www.dondeteesperan.es}, y cuenta con las siguientes características:

En primer lugar, el usuario dispondrá de dos formas para descubrir lugares en base a sus gustos: a través de un cuestionario, y a través de la introducción de una entrada para el sistema de recomendación basado en contenido.

\subsubsection{Cuestionario}

Con la intención de permitir a los usuarios encontrar municipios en base a una búsqueda de criterios que puedan ser de su interés, se diseñan las siguientes 14 cuestiones y filtros, todos ellos opcionales:

\begin{itemize}
    \item Filtrado por comunidad autónoma a la que pertenece el municipio.
    \item Filtrado por provincia a la que pertenece el municipio.
    \item Filtrado por cercanía a la playa del municipio, preguntado en escala de tipo Likert de 5 puntos, donde cada valor se corresponderá con un valor de línea de costa asignado según la heurística utilizada.
    \item Filtrado por cercanía a la montaña del municipio, análogo al anterior.
    \item Filtrado por servicios sanitarios en el municipio, en función de si se desean: centros de salud, centros de salud con servicios de urgencias, hospitales, o todo lo anterior.
    \item Filtrado por servicios educativos en el municipio, en función de si se desean: colegios, universidades o ambos.
    \item Filtrado por cantidad de población del municipio, en función de si el tamaño deseado entra en alguno de los siguientes intervalos: menos de 1000 habitantes, entre 1000 y 5000 habitantes, entre 5000 y 10 000 habitantes, entre 10 000 habitantes y 100 000 habitantes, entre 100 000 habitantes y 250 000 habitantes, y más de 250 000 habitantes.
    \item Filtrado por temperatura media en invierno en el municipio, ofreciendo los siguientes intervalos: menos de 5º C, entre 5 y 10º C, entre 10 y 15º C, y más de 15º C.
    \item Filtrado por temperatura media en verano en el municipio, análogo al anterior, con los siguientes intervalos: menos de 20º C, entre 20 y 22º C, entre 22 y 25º C, y más de 25º C.
    \item Filtrado por altitud del municipio, incluyendo los siguientes intervalos: menos de 250 metros, entre 250 y 500 metros, entre 500 y 750 metros, y más de 750 metros.
    \item Filtrado por distancia a la capital de provincia, con los intervalos: menos de 5 kilómetros, entre 5 y 15 kilómetros, entre 15 y 25 kilómetros, entre 25 y 35 kilómetros, entre 35 y 50 kilómetros, y más de 50 kilómetros.
    \item Filtrado por renta per cápita anual del municipio, dando los intervalos: menos de 20 000 €, entre 20 000 y 25 000 €, entre 25 000 y 30 000 €, y más de 30 000 €.
    \item Filtrado por cobertura de Internet disponible en el municipio, con las opciones: Cobertura móvil 3G en más de la mitad del territorio, cobertura móvil 4G en más de la mitad del territorio, cobertura de fibra óptica de 30 MB en más la mitad del territorio, cobertura de fibra óptica de 100 MB más de la mitad del territorio o cobertura móvil 4G y fibra óptica de 100 MB más de la mitad del territorio
    \item Filtrado por preferencia del usuario entre naturaleza o servicios comerciales, que filtra por los lugares de tipo \guillemotleft sitios\_actividad\guillemotright, \guillemotleft sitios\_ocio\guillemotright\space y \guillemotleft sitios\_natural\guillemotright\space para la primera opción, y por municipios que tengan lugares de tipo \guillemotleft sitios\_comercio\guillemotright, \guillemotleft sitios\_servicio\guillemotright\space y \guillemotleft sitios\_catering\guillemotright\space para la segunda.
\end{itemize}

Además, se consideró necesario añadir una página de \guillemotleft Metodología\guillemotright\space al sitio Web, que explica cómo se ha realizado este trabajo, y que incluye las siguientes notas en relación con el cuestionario anterior –además de asuntos ya descritos en esta memoria, como el cálculo de las heurísticas para las distancias o la toma de valores climáticos promedio en febrero y julio para representar el invierno y el verano, respectivamente–:

\textit{En todos los rangos del cuestionario se incluye el/los extremo/s del intervalo correspondiente, por ejemplo: \guillemotleft Municipios de menos de 10 000 habitantes\guillemotright, incluye a los municipios que tienen 10 000 habitantes, de forma análoga con el ejemplo acotado inferiormente (\guillemotleft Municipios de más de 10~000 habitantes\guillemotright)}.

\textit{La denominación \guillemotleft colegios\guillemotright\space incluye indistintamente a \guillemotleft colegios e institutos\guillemotright\space (Centros de educación primaria y secundaria), dado que los datos del Ministerio de Educación no realizan esta distinción.}

\subsubsection{Programación de la aplicación}

Dentro de la creación de la página Web, se destacan como aspectos relevantes los siguientes:

\begin{enumerate}
    \item Al igual que en la fase ETL (extracción, transformación y limpieza), en la aplicación se hace uso de archivos no añadidos al repositorio público, que mantienen las variables de entorno con credenciales y otros secretos como medida de seguridad.
    \item Se especifica la configuración de la instancia de Google App Engine en la que se ejecutará la aplicación.
    \item Se mantiene el modelo y la base de datos con la que trabajará la aplicación, en formatos CSV y JSON por conveniencia según el procesamiento a realizar. Es importante destacar que ambos archivos solo son accesibles por el servidor, y no se exponen al público para evitar una fuga de datos.
    \item La API maneja únicamente el identificador del municipio para realizar los procesamientos necesarios (cálculo de similitud, búsqueda con filtros aplicados, búsqueda de sugerencias con los sistemas de recomendación, obtención de sugerencias aleatorias...), y es en el último momento, cuando se ha determinado qué municipio se presentará, cuando se resuelve la información y se extrae la necesaria de la base de datos. De esta manera se consigue minimizar la latencia y mejorar la experiencia de uso.
    \item Como se explica posteriormente, debido a las limitaciones para operar con archivos de texto de una Plataforma como Servicio como Google App Engine, se ha programado una API REST sencilla que realiza operaciones de lectura y escritura en archivos CSV en un servidor tradicional de pila LAMP. De esta forma, se conserva la posibilidad de tener archivos de fácil manejo por los sistemas de recomendación, a la par que se evita el acoplamiento con los servicios de la plataforma de computación en la nube elegida, y se reduce el coste de operación de la misma.
    \item Se realizan operaciones de homogeneización de los nombres de los municipios, de forma que la página de información de cada uno pueda estar en un directorio estático conocido y permanente.
    \item Como estas páginas son estáticas, se incluye un método para regenerarlas cuando se necesite, así como otro para probar los cambios en un municipio de ejemplo, sin necesidad de regenerar las miles de páginas estáticas de cada municipio. Como son parte de la programación y depuración, estos métodos se han configurado para que solo funcionen en el entorno local.
    \item Las páginas de los municipios, así como cualquier otra página que reciba información dinámica, se han construido mediante plantillas de Jinja2~\cite{jinja2}, que favorecen la reutilización y escalabilidad mediante componentes de código.
    
\end{enumerate}

Una descripción exhaustiva de cada archivo que compone el directorio Web, y una explicación más detallada de los aspectos anteriores pueden encontrarse en el Anexo \ref{programador} \guillemotleft Documentación técnica de programación\guillemotright. 

\subsubsection{Infraestructura}

Muchos de los retos del sistema han venido dados por el uso de una plataforma como servicio para desplegar la aplicación para usuarios finales. Si bien esto fue elegido por aportar muchas ventajas (como la ausencia de la gestión de la infraestructura, o el pago de los recursos únicamente usados y no por adelantado o por tarifas con cuotas mínimas), es cierto que a cambio se pierde control de configuración y de acceso al sistema.

En particular, como una plataforma como servicio abstrae la infraestructura subyacente, no es posible conseguir que los archivos de Python lean y escriban en ficheros de texto, como sí lo hacían en el entorno local. La alternativa es usar otras APIs de Google para la persistencia de datos, algo que se descartó por desviarse del objetivo fundamental de este trabajo, y por perder la adaptabilidad que brinda usar archivos de texto multipropósito (como texto plano, CSV o JSON) como medios de persistencia de bajo tamaño. Además, no se consideró deseable acoplar la persistencia de datos a la infraestructura de Google, ya que una gran ventaja del \textit{Cloud Computing} es la capacidad de migrar de un proveedor a otro con mínimos cambios en la configuración del despliegue.

Así, se relegó la persistencia al servidor de pila LAMP, del que ya disponía el alumno, y del que tiene pleno control para poder escribir código PHP que lea y actualice los ficheros de texto. De esta manera, este segundo servidor ejerce de base de datos distribuida y es la única fuente de verdad, intercambiando las operaciones a realizar y sus resultados mediante la API HTTP creada a tal fin y usada por la aplicación de Python. Este es uno de los principios de las arquitecturas diseñadas como microservicios, que favorecen el desacoplamiento entre sistemas con operaciones distribuidas de forma transparente para el usuario.

Otro aspecto relevante a destacar tiene que ver con los recursos del sistema. Por defecto, las aplicaciones desplegadas en App Engine se alojan en instancias de tipo \guillemotleft F1\guillemotright, que tienen un límite de 256 MB de memoria RAM y 600 MHz de CPU. Pese a que el funcionamiento era totalmente correcto y rápido en el entorno local, a la hora de desplegarlo a producción (el entorno de App Engine), la aplicación se mostraba lenta y errática, no pudiendo completar en torno al 40 \% de las peticiones. Tras inspeccionar los archivos de registro de errores de App Engine, se llegó a la conclusión de que el rendimiento de Python manejando los ficheros de gran tamaño de la base de datos de municipios era el causante de los problemas, ya que agotaban los recursos de la instancia y provocaban que Google la abortara, resultando en un error HTTP 500, error del servidor.

Este aspecto quedó totalmente subsanado cuando, desde el archivo \guillemotleft app.yaml\guillemotright\space se escaló el sistema, especificando \guillemotleft F4\_1G\guillemotright\space como instancia a ser desplegada, que ofrece 2048 MB de memoria RAM y 2,4 GHz de CPU.

Además, las principales plataformas de \textit{Cloud Computing} abaratan los costes para el desarrollador ofreciendo pago por uso, por lo que cuando el servicio no está siendo accedido se destruye su instancia para que no consuma recursos. La contrapartida de esto es el mayor tiempo que toma la primera petición cuando no hay una instancia creada, lo que también puede impactar en la facturación: como cada instancia está activa unos minutos desde que se crea, y se factura por unidad de tiempo con recursos activos, recibir peticiones periódicas es perjudicial para los costos del sistema. A pesar de ello, como cuando la instancia está destruida la petición lleva un tiempo considerable (varios segundos, afectando a la imagen del producto), se ha aprovechado el segundo servidor de pila LAMP para programar unos trabajos periódicos (\textit{cron jobs}) que realicen unas peticiones a la máquina de App Engine cada hora durante parte de la mañana y parte de la tarde de cada día, a fin de que los usuarios que puedan visitar el sistema en esa franja tengan una mejor experiencia de uso. Dado que el tráfico del sistema es muy bajo, el coste adicional de esta medida se ha considerado despreciable para los beneficios que puede conllevar.

\subsubsection{Aspectos de seguridad técnica y jurídica}

Además de garantizar el correcto funcionamiento de la aplicación Web, se han querido también aplicar los conocimientos de las asignaturas \guillemotleft Fundamentos de Ciberseguridad\guillemotright\space y \guillemotleft Derecho en Seguridad de Datos\guillemotright\space para garantizar la seguridad técnica y jurídica del proyecto. Además de las medidas técnicas señaladas a lo largo de este trabajo, y detalladas posteriormente en el Anexo \ref{programador}, (establecimiento de los permisos mínimos imprescindibles para realizar las operaciones, contraseñas y credenciales seguras no expuestas en el código, uso de claves renovables o sistemas de copias de seguridad automáticos –activados en los dos servidores y disponible en el código a través del control de versiones del repositorio–), se desea comentar también las medidas tomadas desde un punto de vista de la seguridad jurídica.

Es importante señalar que el sitio Web está configurado para redireccionar siempre a la versión HTTPS, cifrando el tráfico intercambiado como medida de seguridad, y convirtiendo las URLs sin \guillemotleft WWW\guillemotright\space a URLs con ellas para facilitar la uniformidad de los recursos y garantizar la seguridad de los usuarios.

Por otra parte, para identificar a los usuarios y sus acciones –estas son, las valoraciones con las que se alimenta el sistema de recomendación colaborativo basado en ítems–, es necesario utilizar una \textit{cookie}. Esto es un pequeño fichero de texto que se almacena en el navegador Web del usuario, y que almacena el identificador único asignado a ese usuario. Sin embargo, dado que el sistema está orientado hacia usuarios europeos, es de aplicación el Reglamento General de Protección de Datos, transpuesto a la legislación española por la Ley Orgánica 3/2018, de 5 de diciembre, de Protección de Datos Personales y Garantía de los Derechos Digitales, que, entre otras cosas, establece la obligación de los prestadores de servicio de recabar el consentimiento informado explícito del usuario antes de realizar operaciones de este tipo~\cite{GDPR_ESP}.

Así, el identificador de usuario –y, por lo tanto, el almacenamiento de la \textit{cookie} en su navegador– solo ocurre cuando este ha dado su consentimiento en una ventana emergente que aparece cuando el sitio carga por primera vez. Esta ventana está gestionada por TermsFeed~\cite{TermsFeed}, y para no volver a mostrarse una vez el usuario ya ha respondido, almacena una \textit{cookie} de tipo técnico con las preferencias manifestadas por el usuario. Solo en el caso de que el usuario otorgue su consentimiento de seguimiento para fines de personalización, se recogerán sus valoraciones y se almacenará la \textit{cookie} propia en su navegador. Además, si el usuario permite el seguimiento para fines analíticos, también se almacenarán \textit{cookies} de Google Analytics~\cite{analytics}, para entender mejor el comportamiento de los usuarios por el sitio Web.

Por su parte, también es necesario disponer de una política de privacidad y \textit{cookies}, que explique transparentemente cómo el sitio Web y su responsable manejan la información de los usuarios. Para ello, se ha contado con el apoyo de la herramienta de Zimrre \& Freelance~\cite{privacidad}.

\subsubsection{Evaluación de la aplicación}

Para evaluar la actividad de despliegue se llevaron a cabo entrevistas con 10 personas de diferente perfil técnico y social, a las que se pidió utilizar el sitio Web sin indicaciones mientras se observa su comportamiento e impresiones. Todos los participantes mostraron su interés en el sistema, y apreciaron su utilidad y facilidad de uso. En particular, el 100 \% de ellos completó las tareas deseadas sin conocer previamente el sistema, y el 85 \% de ellos las completó en menos de 5 minutos cada una.

Además, gracias a su retroalimentación cualitativa y a su observación, se pudieron extraer las siguientes ideas de mejora:

\begin{itemize}
    \item Mejorar el rendimiento del sistema en App Engine, dado que muchas tareas exigían de reintentos por los comentados fallos en el servidor.
    \item Bloquear el desplazamiento lateral en dispositivos táctiles, como tabletas, en el cuestionario para mejorar la navegación.
    \item Reformular preguntas y opciones del cuestionario que presentaron más dudas.
    \item Aclarar el funcionamiento de los controles de desplazamiento para facilitar su uso.
    \item Mejorar la experiencia de usuario en dispositivos móviles.
    \item Modificar las reglas de exclusión de imágenes extraídas automáticamente de Wikipedia para evitar mapas como fondo de página en algunos municipios.
\end{itemize}

Todas ellas fueron incluidas posteriormente en el sistema y se encuentran actualmente publicadas.

\capitulo{6}{Trabajos relacionados}

A continuación, se citarán diversos trabajos relacionados analizados durante la elaboración de este trabajo y cuyas características se han tenido en cuenta a lo largo del proyecto.

\section{Repuéblame}

Creado por Javier y Guido García, Repuéblame~\cite{repueblame} es un portal en línea que ofrece las características de los 2248 municipios de Castilla y León. Entre sus puntos fuertes destacan una interfaz limpia y cuidada y varios selectores para filtrar la información, que parece haber sido recogida de los Datos Abiertos de la Junta de Castilla y León~\cite{datos_abiertos_jcyl}.

\imagen{repueblame2}{Página principal de Repuéblame.}

Destaca también la página de cada municipio, que muestra un mapa y una imagen del lugar, acompañados por un texto generado para cada uno en función de sus características. El servicio otorga a cada localidad una puntuación dentro de un sistema arbitrario no público. Este sistema facilita la lectura y la interpretación de los resultados, pero sesga también la visión del usuario, obligándole a confiar en unos criterios opacos que no tiene por qué compartir.

\imagen{repueblame3}{Página de municipio de Repuéblame.}

El pie de página de cada municipio ofrece otros municipios cercanos y otros similares, utilizando quizás un sistema de recomendación basado en contenido como el descrito en este trabajo. Ofrece también una sección para destacar diferentes municipios de selección editorial con periodicidad mensual.

\imagen{repueblame4}{Pie de página de municipio de Repuéblame.}

Se cree que el sistema podría carecer de \textit{backend}, realizando todo el procesamiento necesario para el filtrado en el cliente Web, lo que ralentiza la navegación y crea problemas de rendimiento en listados grandes. Además, apoya esta hipótesis la existencia de un archivo JavaScript con la base de datos que maneja la página, y que deja entrever las fuentes que se han podido usar y los diferentes criterios seguidos en cada categoría. Este enfoque se considera peligroso ya que permite extraer muy fácilmente el modelo, pieza clave del sistema, haciendo su copia trivial para un tercero malintencionado.

Por último, se destaca que la página de cada municipio está creada estáticamente, lo que mejora el posicionamiento y descubrimiento del sitio, ya que permite encontrar la información de cada municipio desde cualquier buscador.

El trabajo recibió el primer premio de la categoría Productos y Servicios de la V Edición del Concurso de Datos Abiertos de la Junta de Castilla y León~\cite{cyl}.

\section{City Recommender System}

Descrito por Elias Melul en una serie de tres artículos \cite{US1, US2, US3}, se trata de un sistema de recomendación basado en contenido para 303 ciudades estadounidenses. Debido a la concreción de los artículos y a la elaboración del proceso, se ha seguido su propuesta para el desarrollo del sistema presentado, que, a pesar de la gran diferencia en tipo y cantidad de datos con los que trabajan (8131 municipios de muy variados tamaños frente a 303 ciudades), conserva su esencia.

Fruto de la naturaleza de los dos escenarios para los que están pensados ambos sistemas, surgen las principales diferencias. Dado que el recomendador de E. Melul está pensado para ciudades estadounidenses, la existencia de conjuntos de datos de buena calidad y gran cantidad no ha sido problemático. Destacan indicadores socio-económicos, como la seguridad ciudadana, la calidad de vida o el coeficiente de Gini para medir la desigualdad de las localidades. Son todos ellos datos imposibles de conseguir para los municipios españoles, más allá de las ciudades de mayor tamaño, dentro del ámbito de este trabajo; por lo que ha sido necesario considerar cómo se podría sustituir esa información (difícil de entender en su contexto y comparar intuitivamente), por métricas que sí pueden mostrarse junto a cada municipio y que aportan un valor similar.

En el trabajo presentado, los municipios de menor tamaño representan el reto más importante, ya que suponen una mayor dificultad a la hora de encontrar información de múltiples tipos (entradas en Wikipedia, datos tributarios agregados por municipios de similar tamaño, escasas viviendas disponibles...). Este problema, ausente en el trabajo de E. Melul por el diferente dominio, se ha resuelto contando con los datos de las provincias para aquella información que no se ha podido encontrar en un ámbito más pequeño; por ejemplo, a través del número de empleos disponibles en la provincia o el precio medio del alquiler. También se ha encontrado el caso de datos agregados para municipios de muy poco tamaño, como los datos tributarios, donde se ha usado esta información como valor más fidedigno para representar la realidad de la variable de interés.

Se trata de un sistema que carece de interfaz gráfica, ejecutándose únicamente por línea de comandos; por lo que se debió investigar la viabilidad de conectarlo a una página Web antes de comenzar el desarrollo del proyecto presentado. Se encontró el \textit{micro-framework} Flask de Python como mejor opción para realizar esta tarea.

\section{Apadrina un pueblo}

\imagen{apadrina1}{Página principal de Vanwoow.}

Se trata de una iniciativa privada de la empresa de auto-caravanas y otros vehículos de alquiler Vanwoow~\cite{apadrina_un_pueblo}, que ofrece un mapa para localizar pueblos adheridos a la iniciativa, y la posibilidad de sumarse como empresa apadrinadora. Esto último hace que marcas importantes aparezcan como patrocinadoras y colaboradoras del proyecto, del que se presupone una menor recolección automática de datos que en los anteriores debido a la aparente curación de los textos y el contenido mostrado.

\imagen{apadrina2}{Mapa de Vanwoow.}

En el mapa interactivo se puede filtrar por áreas de pernocta, pueblos, experiencias y rutas; y la página de cada municipio, estática de nuevo para ser indexada por los buscadores, ofrece un listado con los principales servicios que el usuario puede encontrarse allí, centrándose en el escenario turístico de pernocta en áreas habilitadas para ello.

\imagen{apadrina3}{Página de municipio en Vanwoow.}

Además, ofrece actividades a realizar en cada municipio y su precio, negocios locales y una sección de comentarios que incentiva la interacción social de la comunidad.

\imagen{apadrina4}{Pasos para anunciar un pueblo en Vanwoow}

Ha recibido varios premios por el carácter innovador de su idea, aunque debido a ser una empresa privada cuyo modelo de negocio gira en torno a la propuesta de descubrir \guillemotleft micropueblos\guillemotright\space en \guillemotleft autocaravana, caravana o camper\guillemotright, se aprecia la curación de contenidos en artículos de blog, en la elaboración de rutas o en la colaboración con los pueblos anunciados, cuyo servicio de promoción se anuncia como una opción para revitalizarlos. Es por ello que se cree que el uso de técnicas de extracción automática de información es limitado y menor que en el caso de otros trabajos estudiados.

\imagen{apadrina5}{Propuesta de valor para municipios de Vanwoow}

\section{Recomendadores de filtro colaborativo}

En cuanto a los recomendadores de filtro colaborativo, se han estudiado los trabajos de Abhinav Ajitsaria~\cite{filtro_colaborativo_1} y Jeffery Chiang~\cite{filtro_colaborativo_2}, que presentan enfoques complementarios a una misma tarea.

El primero apuesta por un filtrado colaborativo basado en usuarios, mientras que el segundo hace uso de un filtrado colaborativo basado en ítems. Ambos trabajos han sido seguidos para la exploración de este tipo de sistemas de recomendación, aunque destaca el mejor rendimiento del segundo sobre el primero, acorde con lo comentado anteriormente: Los sistemas de recomendación basados en ítems son de menor ayuda para usuarios con gustos muy particulares, pero computan menos operaciones y su velocidad de cálculo es mayor.
\capitulo{7}{Conclusiones y Líneas de trabajo futuras}

\section{Conclusiones finales}

A continuación, se presentan las principales conclusiones derivadas de la elaboración de este trabajo.

En primer lugar, el sistema presentado satisface los ocho objetivos marcados inicialmente, aportando una solución innovadora para obtener recomendaciones y perfiles completos de los 8131 municipios del territorio español. Esta solución se apoya en los sistemas de recomendación basados en contenido y de filtro colaborativo basado en ítems, con lo que la potencial utilidad para los usuarios es grande: Podrán descubrir nuevos lugares en base a sus intereses explícitos, el parecido de un municipio con otro y los intereses de otros usuarios similares. Además, el sistema ha buscado ser escalable, seguro y usable; pilares básicos para el éxito de un sistema que hace uso de Big Data, como se ha aprendido a lo largo de numerosas asignaturas del presente máster. Estas impresiones se derivan también de las evaluaciones cuantitativas y cualitativas llevadas a cabo, donde la satisfacción con el sistema presentado es patente.

A nivel técnico, se desprende una conclusión clara: Para generar un buen sistema de recomendación de filtro colaborativo basado en usuarios es imprescindible tener datos reales del comportamiento de las personas cuyos gustos van a ser modelados. Estos datos deben ser cuanto más numerosos mejor, una de las cualidades del Big Data, ya que será esa cantidad la que permitirá generalizar adecuadamente y encontrar los mejores resultados incluso para usuarios con gustos de nicho o que hayan alimentado poco el sistema. Para ello este trabajo sienta también las bases, recogiendo datos de las valoraciones de los usuarios con el fin de poder utilizarlas para dicho propósito en el futuro.

Por otra parte, se pone de manifiesto el valor de la extracción y procesamiento de datos automáticos, y cómo permiten llevar a cabo comparaciones, visualizaciones de datos o extracción de conocimiento de forma comprensible y abarcable por las personas; algo que sería impensable si no fuera por las técnicas de procesamiento de Big Data existentes.

Por último, se observa cómo las opciones de \textit{Cloud Computing} permiten mucha flexibilidad y constituyen una vía escalable de despliegue de proyectos basados en Big Data, con APIs que permiten tener el control mientras se abstrae la complejidad de la infraestructura subyacente; si bien es cierto que presentan otras desventajas inherentes a esta naturaleza, y cuyos retos han sido también parte del desarrollo de este trabajo.

\section{Líneas futuras}

En cuanto a las líneas de trabajo futuras, se destaca la puesta en producción del sistema, de una manera que sea lo más estable y preparada para el futuro posible. Esto abarca, en primer lugar, la creación de pruebas unitarias y búsqueda de la mayor cobertura de \textit{tests} posible; algo que se puede realizar fácilmente con herramientas como UnitTest~\cite{unittest} y Mock~\cite{mock}, bibliotecas estándar para la automatización de pruebas en Python, y SonarQube~\cite{sonar}, herramienta de análisis estático líder en el mercado con amplia integración con \textit{pipelines} de desarrollo tan conocidos como GitHub, por lo que es ideal para este proyecto.

La idea es incrementar la robustez del sistema confiando en ideas sobradamente probadas en el desarrollo de \textit{software}, como los cuadrantes de testeo ágil (\textit{Agile Testing Quadrants},~\cite{agile_testing_quadrants}) o la pirámide de pruebas; que se basa en la fragilidad y coste de los distintos tipos de pruebas automáticos para distribuir su peso a lo largo del sistema~\cite{testing_pyramid}. Con la biblioteca UnitTest para Python y PHPUnit para PHP~\cite{PHPUnit} es posible crear pruebas unitarias y de integración, mientras que herramientas como Selenium~\cite{selenium}, BrowserStack~\cite{browserstack} o Cypress~\cite{cypress} permiten llegar a niveles más profundos, como pruebas de \textit{snapshot} (capturas de pantalla) o pruebas punto-a-punto (\textit{end-to-end tests}).

Además, dado que la extracción de datos se realiza combinando datos públicos (vía conjuntos de datos abiertos y APIs públicas) y \textit{Web scraping}, puede ser útil utilizar herramientas de monitorización, como VisualPing~\cite{visualping} u otras codificadas de forma personalizada a tal efecto, para ser avisado cuando ocurren cambios significativos que podrían requerir una actualización de los datos o de los \textit{scripts} que los extraen.

Es también importante buscar una forma de automatizar el proceso de extracción. Como se vio en \guillemotleft Arquitecturas Big Data\guillemotright, la creación de un \textit{data lake} constituye una de las fases más importantes y que más impacto puede tener en las operaciones de una organización que trabaja con Big Data; ya que permite extraer valor de múltiples maneras de forma constante. A día de hoy, los \textit{scripts} creados para este trabajo deben ser ejecutados en un orden concreto; y sería muy interesante automatizar este proceso para poder ejecutar actualizaciones periódicas de los datos extraídos vía APIs y \textit{Web scraping}, y contar con opciones para incorporar automáticamente también los conjuntos procedentes de datos abiertos y realizar la transformación necesaria para que puedan ser leídos y manejados por el sistema.

Por otra parte, convendría revisar toda la información extraída de Wikipedia, para verificar la exactitud de las informaciones e imágenes cargadas con cada municipio, y asegurarse de que son recursos relevantes y atractivos. Sin embargo, debido a la cantidad de entradas a revisar (8131 localidades), es una ardua tarea cuyo coste podría ser elevado.

Otro aspecto interesante podría ser aumentar la cantidad de recomendaciones ofrecidas a los usuarios de forma proactiva; por ejemplo, mostrando varios municipios similares al actual si el usuario lo valora positivamente, o varios municipios que le podrían gustar basándose en los gustos de otros usuarios y en los suyos a través del sistema híbrido si valora al actual negativamente, en lugar de mostrar una única recomendación personalizada en cada caso. Para aumentar el atractivo de la funcionalidad, estas recomendaciones adicionales se cargarían con una miniatura del municipio, para lo que, idealmente, cada municipio debería tener una imagen asociada garantizada.

Por último, y dada la importancia del mercado internacional en la búsqueda de municipios españoles donde asentarse~\cite{extranjeros_1, extranjeros_2, extranjeros_3}; la traducción del sitio Web, al menos, al inglés se perfila como una característica imprescindible para poder garantizar la escalabilidad y mejora continua del sistema.


%\renewcommand\chaptername{Anexo}
%\renewcommand\thechapter{\Roman{chapter}}
%\setcounter{chapter}{0}

% Añadir entrada en el índice: Anexos
\appendix
\addcontentsline{toc}{part}{Apéndices}
\part*{Apéndices}

\apendice{Plan de Proyecto Software}
\section{Introducción}

A continuación, se explica el plan seguido para desarrollar este proyecto desde varios puntos de vista. 

\section{Planificación temporal}

El proyecto se ha realizado siguiendo la metodología CRISP-DM y combinándola con un desarrollo ágil, caracterizado por ciclos de desarrollo cortos que permiten la iteración continua y la modificación de requisitos de forma natural, y que busca aportar valor al usuario frecuentemente. Para ello, el desarrollo funcional se ha partido en pedazos denominados \textit{user stories}, que siguen el formato habitual \textit{As a, I want to, so that...} para dejar claro quién es el beneficiario principal de la funcionalidad (idealmente el usuario) y expresar qué debe ser construido en términos estándar y sencillos.

\imagen{issue}{Ficha de desarrollo funcional (\textit{issue}) en GitHub siguiendo el formato de \textit{user story}.}

En cuanto a los ciclos de desarrollo, se ha procurado realizar \textit{sprints} (periodos de trabajo que culminan con un hito significativo) de unas dos semanas de duración aproximada, ya que el autor ha compaginado su trabajo a tiempo completo y no siempre ha sido posible avanzar entre semana.

En particular, los hitos de cada iteración han sido los siguientes:

\begin{itemize}
    \item Primera iteración: Del 10 al 25 de diciembre de 2022. Se marcó el objetivo de extraer todos los datos, tanto los procedentes de APIs públicas, como de conjuntos abiertos y \textit{Web scraping}.
    \item Segunda iteración: Del 26 de diciembre de 2022 al 10 de enero de 2023. Durante este periodo la meta marcada fue crear el modelo y la página Web con la que interactuarán los usuarios.
    \item Tercera iteración: Del 11 de enero al 26 de enero de 2023. El objetivo de esta iteración fue comenzar la documentación y avanzarla significativamente.
    \item Cuarta iteración. Del 27 de enero al 12 de febrero de 2023. El objetivo marcado consistió en terminar todo el trabajo, incluyendo la documentación y la presentación.
\end{itemize}

Cada iteración ha estado caracterizada por una reunión y la compartición continua de los artefactos creados con el tutor, que ha actuado de interesado o \textit{stakeholder} en el marco de metodologías ágiles, facilitando el intercambio de información y promoviendo la mejora continua del producto y de los procesos.

\section{Estudio de viabilidad}

La viabilidad del proyecto se analizó desde diversas perspectivas.

\subsection{Viabilidad técnica}

Se estudiaron a fondo los sistemas de recomendación basados en contenido y con filtro colaborativo, tanto basados en usuarios como en artículos. En particular, fue crucial descubrir las publicaciones de Elias Melul~\cite{US1, US2, US3} con ciudades estadounidenses, ya que confirmaron que era posible realizar el trabajo y permitieron establecer una búsqueda de conjuntos de datos similares para cubrir las carencias en cuanto a la cantidad y calidad de los datos por el diferente dominio del problema.

A continuación, se realizó un análisis de las fuentes de datos que se podían encontrar para cada municipio y que se correspondieran de alguna manera con los conjuntos de datos utilizados por E. Melul. Fue determinante explorar satisfactoriamente la viabilidad de poder ejecutar código Python desde un servicio Web, algo que, tras un poco de investigación, resultó ser muy fácil mediante un \textit{framework} Web basado en Python, como Flask~\cite{flask_so}.

Por último, y dado que se quería desplegar el sistema en una plataforma como servicio para aprovechar las ventajas de tarificación a bajo coste que ofrece el \textit{Cloud Computing}, se exploró la posibilidad de desplegar una aplicación basada en Flask en una opción comercial como es Google App Engine. Cuando se vio que era factible de una forma sencilla~\cite{flask_app_engine}, se completó la fase de evaluación de viabilidad técnica.

\subsection{Viabilidad económica}

Debido a las necesidades del proyecto, la viabilidad económica se consiguió de forma inmediata al considerar que el tráfico que soportaría la aplicación sería mínimo, junto a sus necesidades de espacio y procesamiento. Por ello, y haciendo uso de la calculadora de costes de Google Cloud~\cite{calculadora_app_engine}, se estimó que el coste en euros no debería superar la centena al año (entre 50 y 100~€).

Además, al poder aprovechar un servidor de pila LAMP propiedad del alumno, los costes permanecieron estables mientras se ganó mucha versatilidad y funcionalidad.

\subsection{Viabilidad legal}

Aplicando los conocimientos de las asignaturas de la rama de derecho del grado en Ingeniería Informática y de este máster, fue posible conocer los requisitos legales que debía satisfacer este proyecto.

Se debía garantizar la recolección de datos de los usuarios siempre de forma informada y con su previo consentimiento expreso, y ofrecer unas garantías de transparencia que se expresan habitualmente en las políticas de privacidad.

Además, se debe garantizar que se cumplen los términos de uso y atribución de todas las licencias y bibliotecas de terceros de las que se hace uso en el sistema, tanto si son comerciales como si son de código abierto. En particular, la gran mayoría de bibliotecas usadas están licenciadas bajo la licencia MIT o Apache 2.0, que son de las más permisivas en cuanto a sus obligaciones~\cite{MIT},~\cite{Apache2}. Esencialmente, tan solo imponen el requisito de conservar los avisos de las licencias en todas las copias distribuidas del software, no considerando a tal efecto la distribución online de una aplicación Web. Esta es la principal diferencia con AGPL 3.0, una de las licencias más restrictivas que, como parte de la familia de licencias GPL (\textit{GNU General Public License}), exige que el software que las usa se distribuya con la misma licencia (incluyendo la liberación del código fuente, algo referido a veces como \guillemotleft el caracter vírico de la GPL\guillemotright\space por su impacto nocivo en software comercial~\cite{licencia_virica}, ~\cite{GPL}), y además considera el software online como software distribuido~\cite{AGPL}.

Por otro lado, el proyecto hace uso de extracción automática de datos, con lo que era importante confirmar que se podía hacer. Si bien es cierto que hay ciertas lagunas, parece que siempre que no medie interés comercial o que pueda suponer un hecho de competencia desleal, la extracción automática es posible~\cite{scraping_legal}. Como en este caso se trata de un proyecto académico sin fines comerciales, y que no busca establecer ninguna relación de competencia frente a ningún actor cuyos datos se hayan usado, parecía claro que el proyecto era viable legalmente y se dio término a esta fase.
\apendice{Especificación de Requisitos}


\apendice{Especificación de diseño}


\apendice{Documentación técnica de programación}

\section{Introducción}

En el presente anexo se detalla la información necesaria para mantener el sistema por otros programadores.

\section{Estructura de directorios}

Tras clonar el repositorio, se pueden encontrar las carpetas \texttt{Scripts}, \texttt{Data} y \texttt{Web}.

La primera contiene, dentro del directorio \texttt{ETL} (siglas de \texttt{Extracción, transformación y limpieza}) los archivos de Python y PHP con los que se realiza la extracción de datos vía APIs y \textit{Web scraping} y las operaciones de tratamiento de los conjuntos de datos descargados directamente. A su vez, se encuentran también en el directorio \texttt{Scripts} los programas para crear y normalizar el modelo.

Los resultados de las extracciones y procesamientos se encuentran en formato CSV en la carpeta \texttt{Data}. Dentro de esta carpeta están los directorios \texttt{others} y \texttt{support}. En el primero se encuentra la base de datos (el modelo no normalizado) y la base de datos de información extraída de Wikipedia, ambos en formato JSON. Por otra parte, en el segundo directorio se encuentran archivos CSV intermedios o archivos en formato XLSX procedentes de diversos conjuntos de datos.

El directorio \texttt{Web} contiene todo lo relacionado con el producto desplegado. Destacan las carpetas \texttt{api}, con los archivos usados en el servidor secundario de pila LAMP; \texttt{templates}, con las plantillas de Jinja2 usadas por Flask; y \texttt{static}, que contiene las páginas estáticas del sistema y los archivos necesarios para interpretarlas en el navegador. Esto incluye las páginas de cada municipio, presentes en la subcarpeta de su provincia dentro de la carpeta \texttt{municipios}, los archivos de estilo CSS y de JavaScript necesarios para la interacción enriquecida con el sitio Web y las imágenes mostradas alojadas directamente en el servidor.

\section{Manual del programador}

Para ejecutar el servidor local se necesita tener instalado en Python 3 y contar con 2 Gigabytes de espacio libre en disco y en memoria RAM.

Todos los \textit{scripts} están acompañados de un texto que indica cómo usarlos, y que se muestra si se ejecuta el archivo sin especificar ningún parámetro adicional. Además, es posible encontrar comentarios y nombres descriptivos que ilustran el cometido del archivo.

Por razones de universalidad en el acceso al conocimiento, pragmatismo y eficiencia, y con la intención de que el código pueda ser leído, entendido, mantenido o contribuido por personas de cualquier parte del mundo, todos los archivos de programación están escritos en inglés, al igual que las fichas que describen las funciones del sistema (\textit{issues}) en GitHub.

\section{Compilación, instalación y ejecución del proyecto}

Para ejecutar el proyecto localmente se ejecutarán los siguientes comandos desde el directorio \texttt{Web}:

\begin{verbatim}
pip install -r requirements.txt
python3 main.py mode=debug
\end{verbatim}

El primer comando usa el gestor de paquetes Pip de Python para instalar todas las dependencias necesarias del sistema; mientras que el segundo inicia la aplicación localmente (modo \texttt{debug}). Esto es necesario ya que, de lo contrario, la aplicación estaría apuntando hacia el dominio real, por lo que, en realidad, sería equivalente a acceder al sitio en Internet. Para usar la aplicación localmente tan solo se debe acceder a la URL de Flask desde el dispositivo en el que se han ejecutado los comandos anteriores: localhost:5000.

Ejecutar \texttt{python3 main.py mode=production} es equivalente a ejecutar \texttt{python3 main.py}, que es la orden ejecutada automáticamente por App Engine durante el despliegue del sistema.

Las páginas de los municipios se generan accediendo, únicamente desde el modo \texttt{debug}, a localhost:5000/generate-municipalities-pages. En el ordenador del autor, un MacBook Pro de 13 pulgadas de 2016 con macOS 12.6 Monterey, un procesador de 3,1 GHz Intel Core i5 de doble núcleo, 16 Gigabytes de memoria RAM LPDDR3 a 2133 MHz y una tarjeta gráfica Intel Iris Graphics 550 1536 MB, este proceso lleva aproximadamente una hora y media.

Para desplegar el sistema en App Engine tan solo es necesario ejecutar el GitHub Action \texttt{Deploy to GAE}, como se muestra a continuación:

\imagen{deploy}{GitHub Action de despliegue a Google App Engine.}

\section{Pruebas del sistema}

Las pruebas manuales del sistema se pueden realizar accediendo al sitio en \texttt{localhost} e interactuando con él. Se puede obtener una vista de prueba dinámica de la página de municipio –sin necesidad de recrear todas las páginas de forma estática– entrando, únicamente desde el modo \texttt{debug}, a localhost:5000/test-municipality-page.
\apendice{Documentación de usuario}

\section{Introducción}

En el presente anexo se detalla la información relevante para el manejo del sistema por el usuario.

\section{Requisitos de usuarios}

En principio, los usuarios pueden interactuar con el sistema a través de una dirección de Internet, por lo que los requisitos de su dispositivo son mínimos: La aplicación funciona correctamente en ordenadores y dispositivos móviles, aunque se aconseja su uso en dispositivos con pantalla igual o superior a 8 pulgadas para una mejor experiencia.

Si desean ejecutar la aplicación en modo local, los usuarios necesitarán tener instalado en su sistema Python 3 y contar con 2 Gigabytes de espacio libre en disco y en memoria RAM.

\section{Instalación}

Si desean usar el sistema sin necesidad de instalación, tan solo tienen que entrar en el dominio habilitado para ello: www.dondeteesperan.es. Es importante señalar que el sitio Web está configurado para redireccionar las URL sin \guillemotleft www\guillemotright\space a sus equivalentes con el prefijo para facilitar la uniformidad de los recursos y la comodidad de los usuarios.

Si desean usar el sistema localmente, deben clonar el repositorio disponible en github.com/mvarona/TFM\_MUINBDES, abrir un ventana de línea de comandos, colocarse en la carpeta \guillemotleft Web\guillemotright\space y ejecutar los siguientes comandos:

\begin{verbatim}
pip install -r requirements.txt
python3 main.py mode=debug
\end{verbatim}

El primer comando usa el gestor de paquetes Pip de Python para instalar todas las dependencias necesarias del sistema; mientras que el segundo inicia la aplicación localmente (modo \texttt{debug}). Esto es necesario ya que, de lo contrario, la aplicación estaría apuntando hacia el dominio real, por lo que, en realidad, sería equivalente a acceder al sitio en Internet. Para usar la aplicación localmente tan solo se debe acceder a la URL de Flask desde el dispositivo en el que se han ejecutado los comandos anteriores: localhost:5000.

\section{Manual del usuario}

\subsection{Página principal}

Una vez el usuario accede al sistema puede ver la página principal, que muestra una imagen aleatoria de entre los cinco fondos aleatorios. En la página principal se destacan tres funciones principales del sistema: búsqueda según los gustos del usuario de acuerdo con sus respuestas a un formulario, búsqueda según el parecido con otro municipio, y municipio aleatorio o sorpresa.

En el pie de todas las páginas del sitio se encuentra la referencia de la fotografía utilizada en la vista actual y el enlace a la página donde se explica la metodología usada.

Desde cualquier página del sitio Web se puede volver a la principal pulsando en el nombre del dominio de la esquina superior izquierda.

A través de la pancarta emergente inferior para la gestión de las \textit{cookies} es posible acceder a la política de privacidad y \textit{cookies} del sitio Web. Una vez se haya manifestado el consentimiento mediante los botones de aceptar o renunciar, es posible acceder a dicha política, así como al detalle de las \textit{cookies} usadas, desde dondeteesperan.es/privacidad.

\imagen{home}{Página principal.}

\subsection{Cuestionario}

Desde la función \guillemotleft Lugar solo para mí\guillemotright\space el usuario puede filtrar por diversos criterios a través de catorce preguntas de diversa índole y forma de respuesta. En algunas preguntas será necesario elegir una respuesta de un selector vertical pulsando sobre él y después sobre la opción deseada. En otras, se deberá ajustar el selector horizontal al nivel deseado; para ello, se debe presionar en el indicador y arrastrar hacia un lado u otro para soltar sobre el nivel elegido. La apariencia o comportamiento específico de estos controles varía según el sistema operativo, el navegador u otros ajustes del usuario.

\imagen{survey1}{Ejemplo de pregunta mediante selector vertical del cuestionario.}

El usuario se puede desplazar a lo largo de las preguntas mediante las flechas que aparecen en el lado izquierdo y derecho del formulario. El desplazamiento es circular, por lo que después de la última pregunta se pasará a la primera. Unos indicadores rectangulares muestran en qué pregunta se encuentra el usuario con respecto al total, destacando la posición de la pregunta actual en gris y dejando los demás en blanco claro.

\imagen{survey2}{Ejemplo de pregunta mediante selector deslizante del cuestionario.}

Si no se desea filtrar por ninguna opción basta con pasar todas las preguntas sin utilizar ningún selector y dejando las opciones de los controles por defecto.

\imagen{survey3}{Pregunta final del cuestionario.}

En la última pregunta se pulsará en \guillemotleft ¡Vamos!\guillemotright\space para obtener el resultado del filtrado. A continuación, aparecerá la pantalla de carga del resultado sugerido.

\imagen{survey4}{Pantalla de carga del cuestionario.}

Si ha sido posible encontrar un subconjunto resultante de aplicar todas las opciones elegidas anteriormente se mostrará uno de los municipios de dicho grupo. En caso contrario, se mostrará un mensaje al usuario indicando que su selección contiene criterios que excluyen a todos los municipios, por lo que debe ajustarlos de nuevo para poder recomendarle municipios en base a sus criterios.

\imagen{survey5}{Pantalla de ausencia de resultados al aplicar los filtros.}

\subsection{Búsqueda de municipios similares}

A través de la función \guillemotleft Lugar parecido a otro\guillemotright\space el usuario puede obtener el municipio más parecido a uno dado. Para ello, solo tiene que empezar a introducir el nombre del municipio y su provincia, separados por coma, en la caja de texto que aparece en la página. Para elegir el municipio, se debe pulsar en el municipio deseado de los que aparecen en el desplegable de auto-completado, pensado para eliminar errores y facilitar la experiencia de usuario. Si el resultao deseado no aparece, tan solo escriba algunos caracteres más del nombre y provincia del municipio.

Una vez se pulse en la localidad deseada, se mostrará la pantalla de carga y, tras ella, será redirigido a la página del municipio resultante.

\imagen{similar1}{Pantalla para introducir el municipio en base al que se desea obtener la recomendación.}

\subsection{Página de municipio}

La página de municipio contiene la principal información para hacerse una idea del municipio sugerido. Siempre que sea posible, la página cargará una de las imágenes de Wikipedia como fondo de la página. En caso contrario, se cargará una de las imágenes aleatorias del sitio Web.

\imagen{similar2}{Primera mitad de la página de municipio.}

En la primera mitad superior izquierda se observa, si es posible, un extracto de la página del municipio en Wikipedia, junto con un enlace a dicho artículo y a la búsqueda del municipio en Google. Justo debajo se puede encontrar una lista de las etiquetas extraídas automáticamente para ese municipio.

A continuación, se encuentra el panel de valoración del municipio, que permite expresar una valoración positiva o negativa para el municipio mostrado. Es importante señalar que la valoración solo irá asociada al identificador del usuario (generado aleatoriamente en base a su navegador) en caso de que se hayan aceptado las \textit{cookies} para fines estadísticos.

Tras pulsar en \guillemotleft Me gusta\guillemotright, se obtiene una recomendación personalizada para el usuario en base al sistema de recomendación híbrida que agrupa resultados del sistema de recomendación basado en contenido y del sistema de recomendación de filtro colaborativo basado en ítems. Dicha recomendación viene en forma de enlace que, cuando se pulsa, carga la página del municipio sugerido.

\imagen{municipality2-1}{Recomendación ofrecida tras valoración positiva de un municipio sugerido.}

Los dos paneles contiguos inferiores señalan la ubicación del municipio en un mapa de España y muestran una vista del mismo a pie de calle mediante Google Maps y Google Street View.

\imagen{similar3}{Segunda mitad de la página de municipio.}

Paralelamente, en la parte derecha se dispone las principales características del municipio, como datos relacionados con el mercado inmobilario y de trabajo, los servicios sanitarios y educativos, el clima, la altitud, la distancia a la capital de provincia, la renta anual bruta \textit{per cápita}, las horas de sol promedio a lo largo del año o la conectividad a Internet de la que dispone.

A continuación, se muestra una galería de fotos extraídas de su página en Wikipedia, si dispone de ellas.

La página de municipio está especialmente optimizada para dispositivos móviles, donde ajusta la disposición de elementos para fusionar las columnas y obtener una navegación vertical más cómoda.

\imagen{mobile1}{Ejemplo de la primera mitad de la página de municipio adaptada a pantallas móviles.}
\imagen{mobile2}{Ejemplo de la segunda mitad de la página de municipio adaptada a pantallas móviles.}

\subsection{Municipio sorpresa}

La función \guillemotleft Lugar sorpresa\guillemotright\space carga un municipio aleatorio de la base de datos, yendo directamente a su página de municipio.

\subsection{Metodología}

La página de \guillemotleft Metodología\guillemotright\space explica las principales características del producto publicado, haciendo hincapié en que se trata del resultado de este trabajo Fin de Máster, indicando las fuentes de los datos, el proceso llevado a cabo y otra información útil, como diversos criterios de las respuestas del formulario o que la extracción de datos se ha realizado de forma automática no supervisada. Además, se incluyen diversas formas para contactar con el autor para poder obtener más información o reportar algún error.

\imagen{metodologia}{Página de descripción de la metodología usada.}


\bibliographystyle{unsrt}
\bibliography{bibliografia}

\end{document}

% Las referencias se incluyen en el texto usando cite \cite{wiki:latex}. Para citar webs, artículos o libros \cite{koza92}.
% \imagen{escudoInfor}{Autómata para una expresión vacía}
